
\chapter{Примеры написания}\label{ch:ch5}
Мы можем сделать \textbf{жирный текст} и \textit{курсив}.

Сошлёмся на библиографию.
Одна ссылка: \cite[с.~54]{Sokolov}\cite[с.~36]{Gaidaenko}.
Две ссылки: \cite{Sokolov,Gaidaenko}.
Ссылка на собственные работы: \cite{vakbib1, confbib2}.
Много ссылок: %\cite[с.~54]{Lermontov,Management,Borozda} % такой «фокус»
%вызывает biblatex warning относительно опции sortcites, потому что неясно, к
%какому источнику относится уточнение о страницах, а bibtex об этой проблеме
%даже не предупреждает
\cite{Lermontov, Management, Borozda, Marketing, Constitution, FamilyCode,
    Gost.7.0.53, Razumovski, Lagkueva, Pokrovski, Methodology, Berestova,
    Kriger}%
\ifnumequal{\value{bibliosel}}{0}{% Примеры для bibtex8
    \cite{Sirotko, Lukina, Encyclopedia, Nasirova}%
}{% Примеры для biblatex через движок biber
    \cite{Sirotko2, Lukina2, Encyclopedia2, Nasirova2}%
}%
.
И~ещё немного ссылок:~\cite{Article,Book,Booklet,Conference,Inbook,Incollection,Manual,Mastersthesis,
    Misc,Phdthesis,Proceedings,Techreport,Unpublished}
% Следует обратить внимание, что пробел после запятой внутри \cite{}
% обрабатывается ожидаемо, а пробел перед запятой, может вызывать проблемы при
% обработке ссылок.
\cite{medvedev2006jelektronnye, CEAT:CEAT581, doi:10.1080/01932691.2010.513279,
    Gosele1999161,Li2007StressAnalysis, Shoji199895, test:eisner-sample,
    test:eisner-sample-shorted, AB_patent_Pomerantz_1968, iofis_patent1960}%
\ifnumequal{\value{bibliosel}}{0}{% Примеры для bibtex8
}{% Примеры для biblatex через движок biber
    \cite{patent2h, patent3h, patent2}%
}%
.

\ifnumequal{\value{bibliosel}}{0}{% Примеры для bibtex8
Попытка реализовать несколько ссылок на конкретные страницы
для \texttt{bibtex} реализации библиографии:
[\citenum{Sokolov}, с.~54; \citenum{Gaidaenko}, с.~36].
}{% Примеры для biblatex через движок biber
Несколько источников (мультицитата):
% Тут специально написано по-разному тире, для демонстрации, что
% применение специальных тире в настоящий момент в biblatex приводит к непоказу
% "с.".
\cites[vii--x, 5, 7]{Sokolov}[v"--~x, 25, 526]{Gaidaenko}[vii--x, 5, 7]{Techreport},
работает только в \texttt{biblatex} реализации библиографии.
}%

Ссылки на собственные работы:~\cite{vakbib1, confbib1}.

Сошлёмся на приложения: Приложение~\cref{app:A}, Приложение~\cref{app:B2}.

Сошлёмся на формулу: формула~\cref{eq:equation1}.

Сошлёмся на изображение: рисунок~\cref{fig:knuth}.

Стандартной практикой является добавление к ссылкам префикса, характеризующего тип элемента.
Это не является строгим требованием, но~позволяет лучше ориентироваться в документах большого размера.
Например, для ссылок на~рисунки используется префикс \textit{fig},
для ссылки на~таблицу "--- \textit{tab}.

В таблице \cref{tab:tab_pref} приложения~\cref{app:B4} приведён список рекомендуемых
к использованию стандартных префиксов.

В некоторых ситуациях возникает необходимость отойти от требований ГОСТ по оформлению ссылок на
литературу.
В таком случае можно воспользоваться дополнительными опциями пакета \verb+biblatex+.

Например, в ссылке на книгу~\cite{sobenin_kdv} использование опции \verb+maxnames=4+ позволяет
вывести имена всех четырёх авторов.
По ГОСТ имена последних трёх авторов опускаются.

Кроме того, часто возникают проблемы с транслитерованными инициалами. Некоторые буквы русского
алфавита по правилам транслитерации записываются двумя буквами латинского алфавита (ю-yu, ё-yo и
т.д.).
Такие инициалы \verb+biblatex+ будет сокращать до одной буквы, что неверно.
Поправить его работу можно использовав опцию \verb+giveninits=false+.
Пример использования этой опции можно видеть в ссылке~\cite{initials}.

%\section{Формулы}\label{sec:ch2/sec3}

Благодаря пакету \textit{icomma}, \LaTeX~одинаково хорошо воспринимает
в~качестве десятичного разделителя и запятую (\(3,1415\)), и точку (\(3.1415\)).

\subsection{Ненумерованные одиночные формулы}\label{subsec:ch5/sec1/sub1}

Вот так может выглядеть формула, которую необходимо вставить в~строку
по~тексту: \(x \approx \sin x\) при \(x \to 0\).

А вот так выглядит ненумерованная отдельностоящая формула c подстрочными
и надстрочными индексами:
\[
    (x_1+x_2)^2 = x_1^2 + 2 x_1 x_2 + x_2^2
\]

Формула с неопределенным интегралом:
\[
    \int f(\alpha+x)=\sum\beta
\]

При использовании дробей формулы могут получаться очень высокие:
\[
    \frac{1}{\sqrt{2}+
        \displaystyle\frac{1}{\sqrt{2}+
            \displaystyle\frac{1}{\sqrt{2}+\cdots}}}
\]

В формулах можно использовать греческие буквы:
%Все \original... команды заранее, ради этого примера, определены в Dissertation\userstyles.tex
\[
    \alpha\beta\gamma\delta\originalepsilon\epsilon\zeta\eta\theta%
    \vartheta\iota\kappa\varkappa\lambda\mu\nu\xi\pi\varpi\rho\varrho%
    \sigma\varsigma\tau\upsilon\originalphi\phi\chi\psi\omega\Gamma\Delta%
    \Theta\Lambda\Xi\Pi\Sigma\Upsilon\Phi\Psi\Omega
\]
\[%https://texfaq.org/FAQ-boldgreek
    \boldsymbol{\alpha\beta\gamma\delta\originalepsilon\epsilon\zeta\eta%
        \theta\vartheta\iota\kappa\varkappa\lambda\mu\nu\xi\pi\varpi\rho%
        \varrho\sigma\varsigma\tau\upsilon\originalphi\phi\chi\psi\omega\Gamma%
        \Delta\Theta\Lambda\Xi\Pi\Sigma\Upsilon\Phi\Psi\Omega}
\]

Для добавления формул можно использовать пары \verb+$+\dots\verb+$+ и \verb+$$+\dots\verb+$$+,
но~они считаются устаревшими.
Лучше использовать их функциональные аналоги \verb+\(+\dots\verb+\)+ и \verb+\[+\dots\verb+\]+.

\subsection{Ненумерованные многострочные формулы}\label{subsec:ch5/sec2/sub2}

Вот так можно написать две формулы, не нумеруя их, чтобы знаки <<равно>> были
строго друг под другом:
\begin{align}
    f_W & =  \min \left( 1, \max \left( 0, \frac{W_{soil} / W_{max}}{W_{crit}} \right)  \right), \nonumber \\
    f_T & =  \min \left( 1, \max \left( 0, \frac{T_s / T_{melt}}{T_{crit}} \right)  \right), \nonumber
\end{align}

Выровнять систему ещё и по переменной \( x \) можно, используя окружение
\verb|alignedat| из пакета \verb|amsmath|. Вот так:
\[
|x| = \left\{
\begin{alignedat}{2}
     &   & x, \quad & \text{eсли } x\geqslant 0 \\
     & - & x, \quad & \text{eсли } x<0
\end{alignedat}
\right.
\]
Здесь первый амперсанд (в исходном \LaTeX\ описании формулы) означает
выравнивание по~левому краю, второй "--- по~\( x \), а~третий "--- по~слову
<<если>>. Команда \verb|\quad| делает большой горизонтальный пробел.

Ещё вариант:
\[
    |x|=
    \begin{cases}
        \phantom{-}x, \text{если } x \geqslant 0 \\
        -x, \text{если } x<0
    \end{cases}
\]

Кроме того, для  нумерованных формул \verb|alignedat| делает вертикальное
выравнивание номера формулы по центру формулы. Например, выравнивание
компонент вектора:
\begin{equation}
    \label{eq:2p3}
    \begin{alignedat}{2}
        {\mathbf{N}}_{o1n}^{(j)} = \,{\sin} \phi\,n\!\left(n+1\right)
        {\sin}\theta\,
        \pi_n\!\left({\cos} \theta\right)
        \frac{
        z_n^{(j)}\!\left( \rho \right)
        }{\rho}\,
         & {\boldsymbol{\hat{\mathrm e}}}_{r}\,+      \\
        +\,
        {\sin} \phi\,
        \tau_n\!\left({\cos} \theta\right)
        \frac{
        \left[\rho z_n^{(j)}\!\left( \rho \right)\right]^{\prime}
        }{\rho}\,
         & {\boldsymbol{\hat{\mathrm e}}}_{\theta}\,+ \\
        +\,
        {\cos} \phi\,
        \pi_n\!\left({\cos} \theta\right)
        \frac{
        \left[\rho z_n^{(j)}\!\left( \rho \right)\right]^{\prime}
        }{\rho}\,
         & {\boldsymbol{\hat{\mathrm e}}}_{\phi}\:.
    \end{alignedat}
\end{equation}

Ещё об отступах. Иногда для лучшей <<читаемости>> формул полезно
немного исправить стандартные интервалы \LaTeX\ с учётом логической
структуры самой формулы. Например в формуле~\cref{eq:2p3} добавлен
небольшой отступ \verb+\,+ между основными сомножителями, ниже
результат применения всех вариантов отступа:
\begin{align*}
    \backslash!             & \quad f(x) = x^2\! +3x\! +2         \\
    \mbox{по-умолчанию}     & \quad f(x) = x^2+3x+2               \\
    \backslash,             & \quad f(x) = x^2\, +3x\, +2         \\
    \backslash{:}           & \quad f(x) = x^2\: +3x\: +2         \\
    \backslash;             & \quad f(x) = x^2\; +3x\; +2         \\
    \backslash \mbox{space} & \quad f(x) = x^2\ +3x\ +2           \\
    \backslash \mbox{quad}  & \quad f(x) = x^2\quad +3x\quad +2   \\
    \backslash \mbox{qquad} & \quad f(x) = x^2\qquad +3x\qquad +2
\end{align*}

Можно использовать разные математические алфавиты:
\begin{align}
    \mathcal{ABCDEFGHIJKLMNOPQRSTUVWXYZ} \nonumber  \\
    \mathfrak{ABCDEFGHIJKLMNOPQRSTUVWXYZ} \nonumber \\
    \mathbb{ABCDEFGHIJKLMNOPQRSTUVWXYZ} \nonumber
\end{align}

Посмотрим на систему уравнений на примере аттрактора Лоренца:

\[
\left\{
\begin{array}{rl}
    \dot x = & \sigma (y-x)  \\
    \dot y = & x (r - z) - y \\
    \dot z = & xy - bz
\end{array}
\right.
\]

А для вёрстки матриц удобно использовать многоточия:
\[
    \left(
        \begin{array}{ccc}
            a_{11} & \ldots & a_{1n} \\
            \vdots & \ddots & \vdots \\
            a_{n1} & \ldots & a_{nn} \\
        \end{array}
    \right)
\]

\subsection{Нумерованные формулы}\label{subsec:ch5/sec3/sub3}

А вот так пишется нумерованная формула:
\begin{equation}
    \label{eq:equation1}
    e = \lim_{n \to \infty} \left( 1+\frac{1}{n} \right) ^n
\end{equation}

Нумерованных формул может быть несколько:
\begin{equation}
    \label{eq:equation2}
    \lim_{n \to \infty} \sum_{k=1}^n \frac{1}{k^2} = \frac{\pi^2}{6}
\end{equation}

Впоследствии на формулы~\cref{eq:equation1, eq:equation2} можно ссылаться.

Сделать так, чтобы номер формулы стоял напротив средней строки, можно,
используя окружение \verb|multlined| (пакет \verb|mathtools|) вместо
\verb|multline| внутри окружения \verb|equation|. Вот так:
\begin{equation} % \tag{S} % tag - вписывает свой текст
    \label{eq:equation3}
    \begin{multlined}
        1+ 2+3+4+5+6+7+\dots + \\
        + 50+51+52+53+54+55+56+57 + \dots + \\
        + 96+97+98+99+100=5050
    \end{multlined}
\end{equation}

Уравнения~\cref{eq:subeq_1,eq:subeq_2} демонстрируют возможности
окружения \verb|subequations| (пакет \verb|amsmath|).
\begin{subequations}
    \label{eq:subeq_1}
    \begin{gather}
        y = x^2 + 1 \label{eq:subeq_1-1} \\
        y = 2 x^2 - x + 1 \label{eq:subeq_1-2}
    \end{gather}
\end{subequations}
Ссылки на отдельные уравнения~\cref{eq:subeq_1-1,eq:subeq_1-2,eq:subeq_2-1}.
\begin{subequations}
    \label{eq:subeq_2}
    \begin{align}
        y & = x^3 + x^2 + x + 1 \label{eq:subeq_2-1} \\
        y & = x^2
    \end{align}
\end{subequations}

\subsection{Форматирование чисел и размерностей величин}\label{sec:units}

Числа форматируются при помощи команды \verb|\num|:
\num{5,3};
\num{2,3e8};
\num{12345,67890};
\num{2,6 d4};
\num{1+-2i};
\num{.3e45};
\num[exponent-base=2]{5 e64};
\num[exponent-base=2,exponent-to-prefix]{5 e64};
\num{1.654 x 2.34 x 3.430}
\num{1 2 x 3 / 4}.
Для написания последовательности чисел можно использовать команды \verb|\numlist| и \verb|\numrange|:
\numlist{10;30;50;70}; \numrange{10}{30}.
Значения углов можно форматировать при помощи команды \verb|\ang|:
\ang{2.67};
\ang{30,3};
\ang{-1;;};
\ang{;-2;};
\ang{;;-3};
\ang{300;10;1}.

Обратите внимание, что ГОСТ запрещает использование знака <<->> для обозначения отрицательных чисел
за исключением формул, таблиц и~рисунков.
Вместо него следует использовать слово <<минус>>.

Размерности можно записывать при помощи команд \verb|\si| и \verb|\SI|:
\si{\farad\squared\lumen\candela};
\si{\joule\per\mole\per\kelvin};
\si[per-mode = symbol-or-fraction]{\joule\per\mole\per\kelvin};
\si{\metre\per\second\squared};
\SI{0.10(5)}{\neper};
\SI{1.2-3i e5}{\joule\per\mole\per\kelvin};
\SIlist{1;2;3;4}{\tesla};
\SIrange{50}{100}{\volt}.
Список единиц измерений приведён в таблицах~\cref{tab:unit:base,
    tab:unit:derived,tab:unit:accepted,tab:unit:physical,tab:unit:other}.
Приставки единиц приведены в~таблице~\cref{tab:unit:prefix}.

С дополнительными опциями форматирования можно ознакомиться в~описании пакета \texttt{siunitx};
изменить или добавить единицы измерений можно в~файле \texttt{siunitx.cfg}.

\begin{table}
    \centering
    \captionsetup{justification=centering} % выравнивание подписи по-центру
    \caption{Основные величины СИ}\label{tab:unit:base}
    \begin{tabular}{llc}
        \toprule
        Название  & Команда          & Символ         \\
        \midrule
        Ампер     & \verb|\ampere|   & \si{\ampere}   \\
        Кандела   & \verb|\candela|  & \si{\candela}  \\
        Кельвин   & \verb|\kelvin|   & \si{\kelvin}   \\
        Килограмм & \verb|\kilogram| & \si{\kilogram} \\
        Метр      & \verb|\metre|    & \si{\metre}    \\
        Моль      & \verb|\mole|     & \si{\mole}     \\
        Секунда   & \verb|\second|   & \si{\second}   \\
        \bottomrule
    \end{tabular}
\end{table}

\begin{table}
    \small
    \centering
    \begin{threeparttable}% выравнивание подписи по границам таблицы
        \caption{Производные единицы СИ}\label{tab:unit:derived}
        \begin{tabular}{llc|llc}
            \toprule
            Название       & Команда               & Символ              & Название & Команда & Символ \\
            \midrule
            Беккерель      & \verb|\becquerel|     & \si{\becquerel}     &
            Ньютон         & \verb|\newton|        & \si{\newton}                                      \\
            Градус Цельсия & \verb|\degreeCelsius| & \si{\degreeCelsius} &
            Ом             & \verb|\ohm|           & \si{\ohm}                                         \\
            Кулон          & \verb|\coulomb|       & \si{\coulomb}       &
            Паскаль        & \verb|\pascal|        & \si{\pascal}                                      \\
            Фарад          & \verb|\farad|         & \si{\farad}         &
            Радиан         & \verb|\radian|        & \si{\radian}                                      \\
            Грей           & \verb|\gray|          & \si{\gray}          &
            Сименс         & \verb|\siemens|       & \si{\siemens}                                     \\
            Герц           & \verb|\hertz|         & \si{\hertz}         &
            Зиверт         & \verb|\sievert|       & \si{\sievert}                                     \\
            Генри          & \verb|\henry|         & \si{\henry}         &
            Стерадиан      & \verb|\steradian|     & \si{\steradian}                                   \\
            Джоуль         & \verb|\joule|         & \si{\joule}         &
            Тесла          & \verb|\tesla|         & \si{\tesla}                                       \\
            Катал          & \verb|\katal|         & \si{\katal}         &
            Вольт          & \verb|\volt|          & \si{\volt}                                        \\
            Люмен          & \verb|\lumen|         & \si{\lumen}         &
            Ватт           & \verb|\watt|          & \si{\watt}                                        \\
            Люкс           & \verb|\lux|           & \si{\lux}           &
            Вебер          & \verb|\weber|         & \si{\weber}                                       \\
            \bottomrule
        \end{tabular}
    \end{threeparttable}
\end{table}

\begin{table}
    \centering
    \begin{threeparttable}% выравнивание подписи по границам таблицы
        \caption{Внесистемные единицы}\label{tab:unit:accepted}

        \begin{tabular}{llc}
            \toprule
            Название        & Команда           & Символ          \\
            \midrule
            День            & \verb|\day|       & \si{\day}       \\
            Градус          & \verb|\degree|    & \si{\degree}    \\
            Гектар          & \verb|\hectare|   & \si{\hectare}   \\
            Час             & \verb|\hour|      & \si{\hour}      \\
            Литр            & \verb|\litre|     & \si{\litre}     \\
            Угловая минута  & \verb|\arcminute| & \si{\arcminute} \\
            Угловая секунда & \verb|\arcsecond| & \si{\arcsecond} \\ %
            Минута          & \verb|\minute|    & \si{\minute}    \\
            Тонна           & \verb|\tonne|     & \si{\tonne}     \\
            \bottomrule
        \end{tabular}
    \end{threeparttable}
\end{table}

\begin{table}
    \centering
    \captionsetup{justification=centering}
    \caption{Внесистемные единицы, получаемые из эксперимента}\label{tab:unit:physical}
    \begin{tabular}{llc}
        \toprule
        Название                & Команда                  & Символ                 \\
        \midrule
        Астрономическая единица & \verb|\astronomicalunit| & \si{\astronomicalunit} \\
        Атомная единица массы   & \verb|\atomicmassunit|   & \si{\atomicmassunit}   \\
        Боровский радиус        & \verb|\bohr|             & \si{\bohr}             \\
        Скорость света          & \verb|\clight|           & \si{\clight}           \\
        Дальтон                 & \verb|\dalton|           & \si{\dalton}           \\
        Масса электрона         & \verb|\electronmass|     & \si{\electronmass}     \\
        Электрон Вольт          & \verb|\electronvolt|     & \si{\electronvolt}     \\
        Элементарный заряд      & \verb|\elementarycharge| & \si{\elementarycharge} \\
        Энергия Хартри          & \verb|\hartree|          & \si{\hartree}          \\
        Постоянная Планка       & \verb|\planckbar|        & \si{\planckbar}        \\
        \bottomrule
    \end{tabular}
\end{table}

\begin{table}
    \centering
    \begin{threeparttable}% выравнивание подписи по границам таблицы
        \caption{Другие внесистемные единицы}\label{tab:unit:other}
        \begin{tabular}{llc}
            \toprule
            Название                  & Команда              & Символ             \\
            \midrule
            Ангстрем                  & \verb|\angstrom|     & \si{\angstrom}     \\
            Бар                       & \verb|\bar|          & \si{\bar}          \\
            Барн                      & \verb|\barn|         & \si{\barn}         \\
            Бел                       & \verb|\bel|          & \si{\bel}          \\
            Децибел                   & \verb|\decibel|      & \si{\decibel}      \\
            Узел                      & \verb|\knot|         & \si{\knot}         \\
            Миллиметр ртутного столба & \verb|\mmHg|         & \si{\mmHg}         \\
            Морская миля              & \verb|\nauticalmile| & \si{\nauticalmile} \\
            Непер                     & \verb|\neper|        & \si{\neper}        \\
            \bottomrule
        \end{tabular}
    \end{threeparttable}
\end{table}

\begin{table}
    \small
    \centering
    \begin{threeparttable}% выравнивание подписи по границам таблицы
        \caption{Приставки СИ}\label{tab:unit:prefix}
        \begin{tabular}{llcc|llcc}
            \toprule
            Приставка & Команда       & Символ      & Степень      &
            Приставка & Команда       & Символ      & Степень        \\
            \midrule
            Иокто     & \verb|\yocto| & \si{\yocto} & \textminus24 &
            Дека      & \verb|\deca|  & \si{\deca}  & 1              \\
            Зепто     & \verb|\zepto| & \si{\zepto} & \textminus21 &
            Гекто     & \verb|\hecto| & \si{\hecto} & 2              \\
            Атто      & \verb|\atto|  & \si{\atto}  & \textminus18 &
            Кило      & \verb|\kilo|  & \si{\kilo}  & 3              \\
            Фемто     & \verb|\femto| & \si{\femto} & \textminus15 &
            Мега      & \verb|\mega|  & \si{\mega}  & 6              \\
            Пико      & \verb|\pico|  & \si{\pico}  & \textminus12 &
            Гига      & \verb|\giga|  & \si{\giga}  & 9              \\
            Нано      & \verb|\nano|  & \si{\nano}  & \textminus9  &
            Терра     & \verb|\tera|  & \si{\tera}  & 12             \\
            Микро     & \verb|\micro| & \si{\micro} & \textminus6  &
            Пета      & \verb|\peta|  & \si{\peta}  & 15             \\
            Милли     & \verb|\milli| & \si{\milli} & \textminus3  &
            Екса      & \verb|\exa|   & \si{\exa}   & 18             \\
            Санти     & \verb|\centi| & \si{\centi} & \textminus2  &
            Зетта     & \verb|\zetta| & \si{\zetta} & 21             \\
            Деци      & \verb|\deci|  & \si{\deci}  & \textminus1  &
            Иотта     & \verb|\yotta| & \si{\yotta} & 24             \\
            \bottomrule
        \end{tabular}
    \end{threeparttable}
\end{table}

\subsection{Заголовки с формулами: \texorpdfstring{\(a^2 + b^2 = c^2\)}{%
        a\texttwosuperior\ + b\texttwosuperior\ = c\texttwosuperior},
    \texorpdfstring{\(\left\vert\textrm{{Im}}\Sigma\left(
            \protect\varepsilon\right)\right\vert\approx const\)}{|ImΣ (ε)| ≈ const},
    \texorpdfstring{\(\sigma_{xx}^{(1)}\)}{σ\_\{xx\}\textasciicircum\{(1)\}}
}\label{subsec:with_math}

Пакет \texttt{hyperref} берёт текст для закладок в pdf-файле из~аргументов
команд типа \verb|\section|, которые могут содержать математические формулы,
а~также изменения цвета текста или шрифта, которые не отображаются в~закладках.
Чтобы использование формул в заголовках не вызывало в~логе компиляции появление
предупреждений типа <<\texttt{Token not allowed in~a~PDF string
    (Unicode):(hyperref) removing...}>>, следует использовать конструкцию
\verb|\texorpdfstring{}{}|, где в~первых фигурных скобках указывается
формула, а~во~вторых "--- запись формулы для закладок.

\section{Рецензирование текста}\label{sec:markup}

В шаблоне для диссертации и автореферата заданы команды рецензирования.
Они видны при компиляции шаблона в режиме черновика или при установке
соответствующей настройки (\verb+showmarkup+) в~файле \verb+common/setup.tex+.

Команда \verb+\todo+ отмечает текст красным цветом.
\todo{Например, так.}

Команда \verb+\note+ позволяет выбрать цвет текста.
\note{Чёрный, } \note[red]{красный, } \note[green]{зелёный, }
\note[blue]{синий.} \note[orange]{Обратите внимание на ширину и расстановку
    формирующихся пробелов, в~результате приведённой записи (зависит также
    от~применяемого компилятора).}

Окружение \verb+commentbox+ также позволяет выбрать цвет.

\begin{commentbox}[red]
    Красный текст.

    Несколько параграфов красного текста.
\end{commentbox}

\begin{commentbox}[blue]
    Синяя формула.

    \begin{equation}
        \alpha + \beta = \gamma
    \end{equation}
\end{commentbox}

\verb+commentbox+ позволяет закомментировать участок кода в~режиме чистовика.
Чтобы убрать кусок кода для всех режимов, можно использовать окружение
\verb+comment+.

\begin{comment}
Этот текст всегда скрыт.
\end{comment}

\section{Работа со списком сокращений и~условных обозначений}\label{sec:acronyms}

С помощью пакета \texttt{nomencl} можно создавать удобный сортированный список
сокращений и условных обозначений во время написания текста. Вызов
\verb+\nomenclature+ добавляет нужный символ или сокращение с~описанием
в~список, который затем печатается вызовом \verb+\printnomenclature+
в~соответствующем разделе.
Для того, чтобы эти операции прошли, потребуется дополнительный вызов
\verb+makeindex -s nomencl.ist -o %.nls %.nlo+ в~командной строке, где вместо
\verb+%+ следует подставить имя главного файла проекта (\verb+dissertation+
для этого шаблона).
Затем потребуется один или два дополнительных вызова компилятора проекта.
\begin{equation}
    \omega = c k,
\end{equation}
где \( \omega \) "--- частота света, \( c \) "--- скорость света, \( k \) "---
модуль волнового вектора.
\nomenclature{\(\omega\)}{частота света\nomrefeq}
\nomenclature{\(c\)}{скорость света\nomrefpage}
\nomenclature{\(k\)}{модуль волнового вектора\nomrefeqpage}
Использование
\begin{verbatim}
\nomenclature{\(\omega\)}{частота света\nomrefeq}
\nomenclature{\(c\)}{скорость света\nomrefpage}
\nomenclature{\(k\)}{модуль волнового вектора\nomrefeqpage}
\end{verbatim}
после уравнения добавит в список условных обозначений три записи.
Ссылки \verb+\nomrefeq+ на последнее уравнение, \verb+\nomrefpage+ "--- на
страницу, \verb+\nomrefeqpage+ "--- сразу на~последнее уравнение и~на~страницу,
можно опускать и~не~использовать.

Группировкой и сортировкой пунктов в списке можно управлять с~помощью указания
дополнительных аргументов к команде \verb+nomenclature+.
Например, при вызове
\begin{verbatim}
\nomenclature[03]{\( \hbar \)}{постоянная Планка}
\nomenclature[01]{\( G \)}{гравитационная постоянная}
\end{verbatim}
\( G \) будет стоять в списке выше, чем \( \hbar \).
Для корректных вертикальных отступов между строками в описании лучше
не~использовать многострочные формулы в~списке обозначений.

\nomenclature{%
    \( \begin{rcases}
        a_n \\
        b_n
    \end{rcases} \)%
}{коэффициенты разложения Ми в дальнем поле соответствующие электрическим и
    магнитным мультиполям}
\nomenclature[a\( e \)]{\( {\boldsymbol{\hat{\mathrm e}}} \)}{единичный вектор}
\nomenclature{\( E_0 \)}{амплитуда падающего поля}
\nomenclature{\( j \)}{тип функции Бесселя}
\nomenclature{\( k \)}{волновой вектор падающей волны}
\nomenclature{%
    \( \begin{rcases}
        a_n \\
        b_n
    \end{rcases} \)%
}{и снова коэффициенты разложения Ми в дальнем поле соответствующие
    электрическим и магнитным мультиполям. Добавлено много текста, так что
    описание группы условных обозначений значительно превысило высоту этой
    группы...}
\nomenclature{\( L \)}{общее число слоёв}
\nomenclature{\( l \)}{номер слоя внутри стратифицированной сферы}
\nomenclature{\( \lambda \)}{длина волны электромагнитного излучения в вакууме}
\nomenclature{\( n \)}{порядок мультиполя}
\nomenclature{%
    \( \begin{rcases}
        {\mathbf{N}}_{e1n}^{(j)} & {\mathbf{N}}_{o1n}^{(j)} \\
        {\mathbf{M}_{o1n}^{(j)}} & {\mathbf{M}_{e1n}^{(j)}}
    \end{rcases} \)%
}{сферические векторные гармоники}
\nomenclature{\( \mu \)}{магнитная проницаемость в вакууме}
\nomenclature{\( r, \theta, \phi \)}{полярные координаты}
\nomenclature{\( \omega \)}{частота падающей волны}

С помощью \verb+nomenclature+ можно включать в~список сокращения,
не~используя их~в~тексте.
% запись сокращения в список происходит командой \nomenclature,
% а не употреблением самого сокращения
\nomenclature{FEM}{finite element method, метод конечных элементов}
\nomenclature{FIT}{finite integration technique, метод конечных интегралов}
\nomenclature{FMM}{fast multipole method, быстрый метод многополюсника}
\nomenclature{FVTD}{finite volume time-domain, метод конечных объёмов
    во~временной области}
\nomenclature{MLFMA}{multilevel fast multipole algorithm, многоуровневый
    быстрый алгоритм многополюсника}
\nomenclature{BEM}{boundary element method, метод граничных элементов}
\nomenclature{CST MWS}{Computer Simulation Technology Microwave Studio
    программа для компьютерного моделирования уравнен Максвелла}
\nomenclature{DDA}{discrete dipole approximation, приближение дискретиных
    диполей}
\nomenclature{FDFD}{finite difference frequency domain, метод конечных
    разностей в~частотной области}
\nomenclature{FDTD}{finite difference time domain, метод конечных разностей
    во~временной области}
\nomenclature{MoM}{method of moments, метод моментов}
\nomenclature{MSTM}{multiple sphere T-Matrix, метод Т-матриц для множества
    сфер}
\nomenclature{PSTD}{pseudospectral time domain method, псевдоспектральный метод
    во~временной области}
\nomenclature{TLM}{transmission line matrix method, метод матриц линий передач}


\section{Одиночное изображение}\label{sec:ch5/sec1}

\begin{figure}[ht]
    \centerfloat{
        \includegraphics[scale=0.27]{latex}
    }
    \caption{TeX.}\label{fig:latex}
\end{figure}

Для выравнивания изображения по-центру используется команда \verb+\centerfloat+, которая является во
многом улучшенной версией встроенной команды \verb+\centering+.

\section{Длинное название параграфа, в котором мы узнаём как сделать две картинки с~общим номером и названием}\label{sec:ch5/sect2}

А это две картинки под общим номером и названием:
\begin{figure}[ht]
    \begin{minipage}[b][][b]{0.49\linewidth}\centering
        \includegraphics[width=0.5\linewidth]{knuth1} \\ а)
    \end{minipage}
    \hfill
    \begin{minipage}[b][][b]{0.49\linewidth}\centering
        \includegraphics[width=0.5\linewidth]{knuth2} \\ б)
    \end{minipage}
    \caption{Очень длинная подпись к изображению,
        на котором представлены две фотографии Дональда Кнута}
    \label{fig:knuth}
\end{figure}

Те~же~две картинки под~общим номером и~названием,
но с автоматизированной нумерацией подрисунков:
\begin{figure}[ht]
    \centerfloat{
        \hfill
        \subcaptionbox[List-of-Figures entry]{Первый подрисунок\label{fig:knuth_2-1}}{%
            \includegraphics[width=0.25\linewidth]{knuth1}}
        \hfill
        \subcaptionbox{\label{fig:knuth_2-2}}{%
            \includegraphics[width=0.25\linewidth]{knuth2}}
        \hfill
        \subcaptionbox{Третий подрисунок, подпись к которому
            не~помещается на~одной строке}{%
            \includegraphics[width=0.3\linewidth]{example-image-c}}
        \hfill
    }
    \legend{Подрисуночный текст, описывающий обозначения, например. Согласно
        ГОСТ 2.105, пункт 4.3.1, располагается перед наименованием рисунка.}
    \caption[Этот текст попадает в названия рисунков в списке рисунков]{Очень
        длинная подпись к второму изображению, на~котором представлены две
        фотографии Дональда Кнута}\label{fig:knuth_2}
\end{figure}

На рисунке~\cref{fig:knuth_2-1} показан Дональд Кнут без головного убора.
На рисунке~\cref{fig:knuth_2}\subcaptionref*{fig:knuth_2-2}
показан Дональд Кнут в головном уборе.

\section{Векторная графика}\label{sec:ch5/vector}

Возможно вставлять векторные картинки, рассчитываемые \LaTeX\ <<на~лету>>
с~их~предварительной компиляцией. Надписи в таких рисунках будут выполнены
тем же~шрифтом, который указан для документа в целом.
На~рисунке~\cref{fig:tikz_example} на~странице~\pageref{fig:tikz_example}
представлен пример схемы, рассчитываемой пакетом \verb|tikz| <<на~лету>>.
Для ускорения компиляции, подобные рисунки могут быть <<кешированы>>, что
определяется настройками в~\verb|common/setup.tex|.
Причём имя предкомпилированного
файла и~папка расположения таких файлов могут быть отдельно заданы,
что удобно, если не~для подготовки диссертации,
то~для подготовки научных публикаций.
\begin{figure}[ht]
    \centerfloat{
        \ifdefmacro{\tikzsetnextfilename}{\tikzsetnextfilename{tikz_example_compiled}}{}% присваиваемое предкомпилированному pdf имя файла (не обязательно)
        \input{Dissertation/images/tikz_scheme.tikz}

    }
    \legend{}
    \caption[Пример \texttt{tikz} схемы]{Пример рисунка, рассчитываемого
        \texttt{tikz}, который может быть предкомпилирован}\label{fig:tikz_example}
\end{figure}

Множество программ имеют либо встроенную возможность экспортировать векторную
графику кодом \verb|tikz|, либо соответствующий пакет расширения.
Например, в GeoGebra есть встроенный экспорт,
для Inkscape есть пакет svg2tikz,
для Python есть пакет tikzplotlib,
для R есть пакет tikzdevice.

\begin{figure}[htbp]
    \centerfloat{
        \ifdefmacro{\tikzsetnextfilename}{\tikzsetnextfilename{pic2}}{}%
        \input{Dissertation/images/scheme.tikz}
    }
    \legend{%
        \textbf{1} "--- кружок с загогулиной;
        \textbf{2} "--- камертоны;
        \textbf{3} "--- кресты;
        \textbf{4} "--- волны;
        \textbf{5} "--- прямоугольники;
        \textbf{5} "--- пронзённый стрелой прямоугольник.%
    }
    \caption{Составная схема \textit{tikz}}\label{fig:scheme-tikz}
\end{figure}

На рисунке~\cref{fig:scheme-tikz} представлена составная схема \textit{tikz}.
Каждый её элемент нарисован в отдельном файле в единичном масштабе.
Расстановка элементов на~рисунке производится при помощи аргументов \texttt{xshift},
\texttt{yshift}, \texttt{rotate} и~\texttt{scale} окружения \texttt{scope}.

Пример использования библиотеки \textit{circuitikz} изображён на рисунке~\cref{fig:circuitikz}.

\begin{figure}[htbp]
    \centerfloat{
        \input{Dissertation/images/circuit.tikz}
    }
    \caption{Схема \textit{circuitikz}}\label{fig:circuitikz}
\end{figure}

Красивые графики также можно добавлять при помощи пакета \textit{pgfplot}~(рисунок~\cref{fig:pgfplot}).
Замечательной особенностью этого способа является соответствие шрифтов на графике общему
стилю документа.

\begin{figure}[htbp]
    \centerfloat{
        \input{Dissertation/images/plot_csv.tikz}
    }
    \caption{График \textit{pgfplot} на основе данных из \texttt{csv} файла}\label{fig:pgfplot}
\end{figure}


\section{Пример вёрстки списков}\label{sec:ch5/sec3}

\noindent Нумерованный список:
\begin{enumerate}
    \item Первый пункт.
    \item Второй пункт.
    \item Третий пункт.
\end{enumerate}

\noindent Маркированный список:
\begin{itemize}
    \item Первый пункт.
    \item Второй пункт.
    \item Третий пункт.
\end{itemize}

\noindent Вложенные списки:
\begin{itemize}
    \item Имеется маркированный список.
          \begin{enumerate}
              \item В нём лежит нумерованный список,
              \item в котором
                    \begin{itemize}
                        \item лежит ещё один маркированный список.
                    \end{itemize}
          \end{enumerate}
\end{itemize}

\noindent Нумерованные вложенные списки:
\begin{enumerate}
    \item Первый пункт.
    \item Второй пункт.
    \item Вообще, по ГОСТ 2.105 первый уровень нумерации
          (при необходимости ссылки в тексте документа на одно из перечислений)
          идёт буквами русского или латинского алфавитов,
          а второй "--- цифрами со~скобками.
          Здесь отходим от ГОСТ.
          \begin{enumerate}
              \item в нём лежит нумерованный список,
              \item в котором
                    \begin{enumerate}
                        \item ещё один нумерованный список,
                        \item третий уровень нумерации не нормирован ГОСТ 2.105;
                        \item обращаем внимание на строчность букв,
                        \item в этом списке
                              \begin{itemize}
                                  \item лежит ещё один маркированный список.
                              \end{itemize}
                    \end{enumerate}

          \end{enumerate}

    \item Четвёртый пункт.
\end{enumerate}

\section{Традиции русского набора}

Много полезных советов приведено в материале
<<\href{https://kostyrka.ru/main/ru/typesetting-and-typography-crash-course-by-kostyrka/}{Краткий курс благородного набора}>>
(автор А.\:В.~Костырка).
Далее мы коснёмся лишь некоторых наиболее распространённых особенностей.

\subsection{Пробелы}

В~русском наборе принято:
\begin{itemize}
    \item единицы измерения, знак процента отделять пробелами от~числа:
          10~кВт, 15~\% (согласно ГОСТ 8.417, раздел 8);
    \item \(\tg 20\text{\textdegree}\), но: 20~{\textdegree}C
          (согласно ГОСТ 8.417, раздел 8);
    \item знак номера, параграфа отделять от~числа: №~5, \S~8;
    \item стандартные сокращения: т.\:е., и~т.\:д., и~т.\:п.;
    \item неразрывные пробелы в~предложениях.
\end{itemize}

\subsection{Математические знаки и символы}

Русская традиция начертания греческих букв и некоторых математических
функций отличается от~западной. Это исправляется серией
\verb|\renewcommand|.
\begin{itemize}
    %Все \original... команды заранее, ради этого примера, определены в Dissertation\userstyles.tex
    \item[До:] \( \originalepsilon \originalge \originalphi\),
          \(\originalphi \originalleq \originalepsilon\),
          \(\originalkappa \in \originalemptyset\),
          \(\originaltan\),
          \(\originalcot\),
          \(\originalcsc\).
    \item[После:] \( \epsilon \ge \phi\),
          \(\phi \leq \epsilon\),
          \(\kappa \in \emptyset\),
          \(\tan\),
          \(\cot\),
          \(\csc\).
\end{itemize}

Кроме того, принято набирать греческие буквы вертикальными, что
решается подключением пакета \verb|upgreek| (см. закомментированный
блок в~\verb|userpackages.tex|) и~аналогичным переопределением в
преамбуле (см.~закомментированный блок в~\verb|userstyles.tex|). В
этом шаблоне такие переопределения уже включены.

Знаки математических операций принято переносить. Пример переноса
в~формуле~\eqref{eq:equation3}.

\subsection{Кавычки}
В английском языке приняты одинарные и двойные кавычки в~виде ‘...’ и~“...”.
В~России приняты французские («...») и~немецкие („...“) кавычки (они называются
«ёлочки» и~«лапки», соответственно). ,,Лапки`` обычно используются внутри
<<ёлочек>>, например, <<... наш гордый ,,Варяг``...>>.

Французкие левые и правые кавычки набираются
как лигатуры \verb|<<| и~\verb|>>|, а~немецкие левые
и правые кавычки набираются как лигатуры \verb|,,| и~\verb|‘‘| (\verb|``|).

Вместо лигатур или команд с~активным символом "\ можно использовать команды
\verb|\glqq| и \verb|\grqq| для набора немецких кавычек и команды \verb|\flqq|
и~\verb|\frqq| для набора французских кавычек. Они определены в пакете
\verb|babel|.

\subsection{Тире}
%  babel+pdflatex по умолчанию, в polyglossia надо включать опцией (и перекомпилировать с удалением временных файлов)
Команда \verb|"---| используется для печати тире в тексте. Оно может быть
несколько короче английского длинного тире (подробности в~документации
русификации babel). Кроме того, команда задаёт небольшую жёсткую отбивку
от~слова, стоящего перед тире. При этом, само тире не~отрывается от~слова.
После тире следует такая же отбивка от текста, как и~перед тире. При наборе
текста между словом и командой, за которым она следует, должен стоять пробел.

В составных словах, таких, как <<Закон Менделеева"--~Клапейрона>>, для печати
тире надо использовать команду \verb|"--~|. Она ставит более короткое,
по~сравнению с~английским, тире и позволяет делать переносы во втором слове.
При~наборе текста команда \verb|"--~| не отделяется пробелом от слова,
за~которым она следует (\verb|Менделеева"--~|). Следующее за командой слово
может быть  отделено от~неё пробелом или перенесено на другую строку.

Если прямая речь начинается с~абзаца, то перед началом её печатается тире
командой \verb|"--*|. Она печатает русское тире и жёсткую отбивку нужной
величины перед текстом.

\subsection{Дефисы и переносы слов}
%  babel+pdflatex по умолчанию, в polyglossia надо включать опцией (и перекомпилировать с удалением временных файлов)
Для печати дефиса в~составных словах введены две команды. Команда~\verb|"~|
печатает дефис и~запрещает делать переносы в~самих словах, а~команда \verb|"=|
печатает дефис, оставляя \TeX ’у право делать переносы в~самих словах.

В отличие от команды \verb|\-|, команда \verb|"-| задаёт место в~слове, где
можно делать перенос, не~запрещая переносы и~в~других местах слова.

Команда \verb|""| задаёт место в~слове, где можно делать перенос, причём дефис
при~переносе в~этом месте не~ставится.

Команда \verb|",| вставляет небольшой пробел после инициалов с~правом переноса
в~фамилии.

\section{Текст из панграмм и формул}

Любя, съешь щипцы, "--- вздохнёт мэр, "--- кайф жгуч. Шеф взъярён тчк щипцы
с~эхом гудбай Жюль. Эй, жлоб! Где туз? Прячь юных съёмщиц в~шкаф. Экс-граф?
Плюш изъят. Бьём чуждый цен хвощ! Эх, чужак! Общий съём цен шляп (юфть) "---
вдрызг! Любя, съешь щипцы, "--- вздохнёт мэр, "--- кайф жгуч. Шеф взъярён тчк
щипцы с~эхом гудбай Жюль. Эй, жлоб! Где туз? Прячь юных съёмщиц в~шкаф.
Экс-граф? Плюш изъят. Бьём чуждый цен хвощ! Эх, чужак! Общий съём цен шляп
(юфть) "--- вдрызг! Любя, съешь щипцы, "--- вздохнёт мэр, "--- кайф жгуч. Шеф
взъярён тчк щипцы с~эхом гудбай Жюль. Эй, жлоб! Где туз? Прячь юных съёмщиц
в~шкаф. Экс-граф? Плюш изъят. Бьём чуждый цен хвощ! Эх, чужак! Общий съём цен
шляп (юфть) "--- вдрызг! Любя, съешь щипцы, "--- вздохнёт мэр, "--- кайф жгуч.
Шеф взъярён тчк щипцы с~эхом гудбай Жюль. Эй, жлоб! Где туз? Прячь юных съёмщиц
в~шкаф. Экс-граф? Плюш изъят. Бьём чуждый цен хвощ! Эх, чужак! Общий съём цен
шляп (юфть) "--- вдрызг! Любя, съешь щипцы, "--- вздохнёт мэр, "--- кайф жгуч.
Шеф взъярён тчк щипцы с~эхом гудбай Жюль. Эй, жлоб! Где туз? Прячь юных съёмщиц
в~шкаф. Экс-граф? Плюш изъят. Бьём чуждый цен хвощ! Эх, чужак! Общий съём цен
шляп (юфть) "--- вдрызг! Любя, съешь щипцы, "--- вздохнёт мэр, "--- кайф жгуч.
Шеф взъярён тчк щипцы с~эхом гудбай Жюль. Эй, жлоб! Где туз? Прячь юных съёмщиц
в~шкаф. Экс-граф? Плюш изъят. Бьём чуждый цен хвощ! Эх, чужак! Общий съём цен
шляп (юфть) "--- вдрызг! Любя, съешь щипцы, "--- вздохнёт мэр, "--- кайф жгуч.
Шеф взъярён тчк щипцы с~эхом гудбай Жюль. Эй, жлоб! Где туз? Прячь юных съёмщиц
в~шкаф. Экс-граф? Плюш изъят. Бьём чуждый цен хвощ! Эх, чужак! Общий съём цен
шляп (юфть) "--- вдрызг! Любя, съешь щипцы, "--- вздохнёт мэр, "--- кайф жгуч.
Шеф взъярён тчк щипцы с~эхом гудбай Жюль. Эй, жлоб! Где туз? Прячь юных съёмщиц
в~шкаф. Экс-граф? Плюш изъят. Бьём чуждый цен хвощ! Эх, чужак! Общий съём цен
шляп (юфть) "--- вдрызг! Любя, съешь щипцы, "--- вздохнёт мэр, "--- кайф жгуч.
Шеф взъярён тчк щипцы с~эхом гудбай Жюль. Эй, жлоб! Где туз? Прячь юных съёмщиц
в~шкаф. Экс-граф? Плюш изъят. Бьём чуждый цен хвощ! Эх, чужак! Общий съём цен
шляп (юфть) "--- вдрызг! Любя, съешь щипцы, "--- вздохнёт мэр, "--- кайф жгуч.
Шеф взъярён тчк щипцы с~эхом гудбай Жюль. Эй, жлоб! Где туз? Прячь юных съёмщиц
в~шкаф. Экс-граф? Плюш изъят. Бьём чуждый цен хвощ! Эх, чужак! Общий съём цен
шляп (юфть) "--- вдрызг! Любя, съешь щипцы, "--- вздохнёт мэр, "--- кайф жгуч.
Шеф взъярён тчк щипцы с~эхом гудбай Жюль. Эй, жлоб! Где туз? Прячь юных съёмщиц
в~шкаф. Экс-граф? Плюш изъят. Бьём чуждый цен хвощ! Эх, чужак! Общий съём цен
шляп (юфть) "--- вдрызг!Любя, съешь щипцы, "--- вздохнёт мэр, "--- кайф жгуч.
Шеф взъярён тчк щипцы с~эхом гудбай Жюль. Эй, жлоб! Где туз? Прячь юных съёмщиц
в~шкаф. Экс-граф? Плюш изъят. Бьём чуждый цен хвощ! Эх, чужак! Общий съём цен

Ку кхоро адолэжкэнс волуптариа хаж, вим граэко ыкчпэтында ты. Граэкы жэмпэр
льюкяльиюч квуй ку, аэквюы продыжщэт хаж нэ. Вим ку магна пырикульа, но квюандо
пожйдонёюм про. Квуй ат рыквюы ёнэрмйщ. Выро аккузата вим нэ.
\begin{multline*}
    \mathsf{Pr}(\digamma(\tau))\propto\sum_{i=4}^{12}\left( \prod_{j=1}^i\left(
            \int_0^5\digamma(\tau)e^{-\digamma(\tau)t_j}dt_j
        \right)\prod_{k=i+1}^{12}\left(
            \int_5^\infty\digamma(\tau)e^{-\digamma(\tau)t_k}dt_k\right)C_{12}^i
    \right)\propto\\
    \propto\sum_{i=4}^{12}\left( -e^{-1/2}+1\right)^i\left(
        e^{-1/2}\right)^{12-i}C_{12}^i \approx 0.7605,\quad
    \forall\tau\neq\overline{\tau}
\end{multline*}
Квуй ыёюз омниюм йн. Экз алёквюам кончюлату квуй, ты альяквюам ёнвидюнт пэр.
Зыд нэ коммодо пробатуж. Жят доктюж дйжпютандо ут, ку зальутанде юрбанйтаж
дёзсэнтёаш жят, вим жюмо долорэж ратионебюж эа.

Ад ентэгры корпора жплэндидэ хаж. Эжт ат факэтэ дычэрунт пэржыкюти. Нэ нам
доминг пэрчёус. Ку квюо ёужто эррэм зючкёпит. Про хабэо альбюкиюс нэ.
\[
    \begin{pmatrix}
        a_{11} & a_{12} & a_{13} \\
        a_{21} & a_{22} & a_{23}
    \end{pmatrix}
\]

\[
    \begin{vmatrix}
        a_{11} & a_{12} & a_{13} \\
        a_{21} & a_{22} & a_{23}
    \end{vmatrix}
\]

\[
    \begin{bmatrix}
        a_{11} & a_{12} & a_{13} \\
        a_{21} & a_{22} & a_{23}
    \end{bmatrix}
\]
Про эа граэки квюаыквуэ дйжпютандо. Ыт вэл тебиквюэ дэфянятйоныс, нам жолюм
квюандо мандамюч эа. Эож пауло лаудым инкедыринт нэ, пэрпэтюа форынчйбюж пэр
эю. Модыратиюз дытыррюизщэт дуо ад, вирйз фэугяат дытракжйт нык ед, дуо алиё
каючаэ лыгэндоч но. Эа мольлиз юрбанйтаж зигнёфэрумквюы эжт.

Про мандамюч кончэтытюр ед. Трётанё прёнкипыз зигнёфэрумквюы вяш ан. Ат хёз
эквюедым щуавятатэ. Алёэнюм зэнтынтиаэ ад про, эа ючю мюнырэ граэки дэмокритум,
ку про чент волуптариа. Ыльит дыкоры аляквюид еюж ыт. Ку рыбюм мюндй ютенам
дуо.
\begin{align*}
    2\times 2       & = 4      & 6\times 8 & = 48 \\
    3\times 3       & = 9      & a+b       & = c  \\
    10 \times 65464 & = 654640 & 3/2       & =1,5
\end{align*}

\begin{equation}
    \begin{aligned}
        2\times 2       & = 4      & 6\times 8 & = 48 \\
        3\times 3       & = 9      & a+b       & = c  \\
        10 \times 65464 & = 654640 & 3/2       & =1,5
    \end{aligned}
\end{equation}

Пэр йн тальэ пожтэа, мыа ед попюльо дэбетиз жкрибэнтур. Йн квуй аппэтырэ
мэнандря, зыд аляквюид хабымуч корпора йн. Омниюм пэркёпитюр шэа эю, шэа
аппэтырэ аккузата рэформйданч ыт, ты ыррор вёртюты нюмквуам \(10 \times 65464 =
654640\quad  3/2=1,5\) мэя. Ипзум эуежмод \(a+b = c\) мальюизчыт ад дуо. Ад
фэюгаят пытынтёюм адвыржаряюм вяш. Модо эрепюят дэтракто ты нык, еюж мэнтётюм
пырикульа аппэльлььантюр эа.

Мэль ты дэлььынётё такематыш. Зэнтынтиаэ конклььюжионэмквуэ ан мэя. Вёжи лебыр
квюаыквуэ квуй нэ, дуо зймюл дэлььиката ку. Ыам ку алиё путынт.

%Большая фигурная скобка только справа
\[\left. %ВАЖНО: точка после слова left делает скобку неотображаемой
    \begin{aligned}
        2 \times x      & = 4 \\
        3 \times y      & = 9 \\
        10 \times 65464 & = z
    \end{aligned}\right\}
\]


Конвынёры витюпырата но нам, тебиквюэ мэнтётюм позтюлант ед про. Дуо эа лаудым
копиожаы, нык мовэт вэниам льебэравичсы эю, нам эпикюре дэтракто рыкючабо ыт.
Вэрйтюж аккюжамюз ты шэа, дэбетиз форынчйбюж жкряпшэрит ыт прё. Ан еюж тымпор
рыфэррэнтур, ючю дольор котёдиэквюэ йн. Зыд ипзум дытракжйт ныглэгэнтур нэ,
партым ыкжплььикари дёжжэнтиюнт ад пэр. Мэль ты кытэрож молыжтйаы, нам но ыррор
жкрипта аппарэат.

\[ \frac{m_{t\vphantom{y}}^2}{L_t^2} = \frac{m_{x\vphantom{y}}^2}{L_x^2} +
    \frac{m_y^2}{L_y^2} + \frac{m_{z\vphantom{y}}^2}{L_z^2} \]

Вэре льаборэж тебиквюэ хаж ут. Ан пауло торквюатоз хаж, нэ пробо фэугяат
такематыш шэа. Мэльёуз пэртинакёа юлламкорпэр прё ад, но мыа рыквюы конкыптам.
Хёз квюот пэртинакёа эи, ельлюд трактатоз пэр ад. Зыд ед анёмал льаборэж
номинави, жят ад конгуы льабятюр. Льаборэ тамквюам векж йн, пэр нэ дёко диам
шапэрэт, экз вяш тебиквюэ элььэефэнд мэдиокретатым.

Нэ про натюм фюйзчыт квюальизквюэ, аэквюы жкаывола мэль ку. Ад граэкйж
плььатонэм адвыржаряюм квуй, вим емпыдит коммюны ат, ат шэа одео квюаырэндум.
Вёртюты ажжынтиор эффикеэнди эож нэ, доминг лаборамюз эи ыам. Чэнзэрет
мныжаркхюм экз эож, ыльит тамквюам факильизиж нык эи. Квуй ан элыктрам
тинкидюнт ентырпрытаряш. Йн янвыняры трактатоз зэнтынтиаэ зыд. Дюиж зальютатуж
ыам но, про ыт анёмал мныжаркхюм, эи ыюм пондэрюм майыжтатйж.

\section{Таблица обыкновенная}\label{sec:ch5/sect1}

Так размещается таблица:

\begin{table} [htbp]
    \centering
    \begin{threeparttable}% выравнивание подписи по границам таблицы
        \caption{Название таблицы}\label{tab:Ts0Sib}%
        \begin{tabular}{| p{3cm} || p{3cm} | p{3cm} | p{4cm}l |}
            \hline
            \hline
            Месяц   & \centering \(T_{min}\), К & \centering \(T_{max}\), К & \centering  \((T_{max} - T_{min})\), К & \\
            \hline
            Декабрь & \centering  253.575       & \centering  257.778       & \centering      4.203                  & \\
            Январь  & \centering  262.431       & \centering  263.214       & \centering      0.783                  & \\
            Февраль & \centering  261.184       & \centering  260.381       & \centering     \(-\)0.803              & \\
            \hline
            \hline
        \end{tabular}
    \end{threeparttable}
\end{table}

\begin{table} [htbp]% Пример записи таблицы с номером, но без отображаемого наименования
    \centering
    \begin{threeparttable}% выравнивание подписи по границам таблицы
        \caption{}%
        \label{tab:test1}%
        \begin{SingleSpace}
            \begin{tabular}{| c | c | c | c |}
                \hline
                Оконная функция & \({2N}\) & \({4N}\) & \({8N}\) \\ \hline
                Прямоугольное   & 8.72     & 8.77     & 8.77     \\ \hline
                Ханна           & 7.96     & 7.93     & 7.93     \\ \hline
                Хэмминга        & 8.72     & 8.77     & 8.77     \\ \hline
                Блэкмана        & 8.72     & 8.77     & 8.77     \\ \hline
            \end{tabular}%
        \end{SingleSpace}
    \end{threeparttable}
\end{table}

Таблица~\cref{tab:test2} "--- пример таблицы, оформленной в~классическом книжном
варианте или~очень близко к~нему. \mbox{ГОСТу} по~сути не~противоречит. Можно
ещё~улучшить представление, с~помощью пакета \verb|siunitx| или~подобного.

\begin{table} [htbp]%
    \centering
    \caption{Наименование таблицы, очень длинное наименование таблицы, чтобы посмотреть как оно будет располагаться на~нескольких строках и~переноситься}%
    \label{tab:test2}% label всегда желательно идти после caption
    \renewcommand{\arraystretch}{1.5}%% Увеличение расстояния между рядами, для улучшения восприятия.
    \begin{SingleSpace}
        \begin{tabular}{@{}@{\extracolsep{20pt}}llll@{}} %Вертикальные полосы не используются принципиально, как и лишние горизонтальные (допускается по ГОСТ 2.105 пункт 4.4.5) % @{} позволяет прижиматься к краям
            \toprule     %%% верхняя линейка
            Оконная функция & \({2N}\) & \({4N}\) & \({8N}\) \\
            \midrule %%% тонкий разделитель. Отделяет названия столбцов. Обязателен по ГОСТ 2.105 пункт 4.4.5
            Прямоугольное   & 8.72     & 8.77     & 8.77     \\
            Ханна           & 7.96     & 7.93     & 7.93     \\
            Хэмминга        & 8.72     & 8.77     & 8.77     \\
            Блэкмана        & 8.72     & 8.77     & 8.77     \\
            \bottomrule %%% нижняя линейка
        \end{tabular}%
    \end{SingleSpace}
\end{table}

\section{Таблица с многострочными ячейками и примечанием}

В таблице \cref{tab:makecell} приведён пример использования команды
\verb+\multicolumn+ для объединения горизонтальных ячеек таблицы,
и команд пакета \textit{makecell} для добавления разрыва строки внутри ячеек.
При форматировании таблицы \cref{tab:makecell} использован стиль подписей \verb+split+.
Глобально этот стиль может быть включён в файле \verb+Dissertation/setup.tex+ для диссертации и в
файле \verb+Synopsis/setup.tex+ для автореферата.
Однако такое оформление не~соответствует ГОСТ.

\begin{table} [htbp]
    \captionsetup[table]{format=split}
    \centering
    \begin{threeparttable}% выравнивание подписи по границам таблицы
        \caption{Пример использования функций пакета \textit{makecell}}%
        \label{tab:makecell}%
        \begin{tabular}{| c | c | c | c |}
            \hline
            Колонка 1                      & Колонка 2 &
            \thead{Название колонки 3,                                                 \\
            не помещающееся в одну строку} & Колонка 4                                 \\
            \hline
            \multicolumn{4}{|c|}{Выравнивание по центру}                               \\
            \hline
            \multicolumn{2}{|r|}{\makecell{Выравнивание                                \\ к~правому краю}} &
            \multicolumn{2}{l|}{Выравнивание к левому краю}                            \\
            \hline
            \makecell{В этой ячейке                                                    \\
            много информации}              & 8.72      & 8.55                   & 8.44 \\
            \cline{3-4}
            А в этой мало                  & 8.22      & \multicolumn{2}{c|}{5}        \\
            \hline
        \end{tabular}%
    \end{threeparttable}
\end{table}

Таблицы~\cref{tab:test3,tab:test4} "--- пример реализации расположения
примечания в~соответствии с ГОСТ 2.105. Каждый вариант со своими достоинствами
и~недостатками. Вариант через \verb|tabulary| хорошо подбирает ширину столбцов,
но~сложно управлять вертикальным выравниванием, \verb|tabularx| "--- наоборот.
\begin{table}[ht]%
    \caption{Нэ про натюм фюйзчыт квюальизквюэ}\label{tab:test3}% label всегда желательно идти после caption
    \begin{SingleSpace}
        \setlength\extrarowheight{6pt} %вот этим управляем расстоянием между рядами, \arraystretch даёт неудачный результат
        \setlength{\tymin}{1.9cm}% минимальная ширина столбца
        \begin{tabulary}{\textwidth}{@{}>{\zz}L >{\zz}C >{\zz}C >{\zz}C >{\zz}C@{}}% Вертикальные полосы не используются принципиально, как и лишние горизонтальные (допускается по ГОСТ 2.105 пункт 4.4.5) % @{} позволяет прижиматься к краям
            \toprule     %%% верхняя линейка
            доминг лаборамюз эи ыам (Общий съём цен шляп (юфть)) & Шеф взъярён &
            адвыржаряюм &
            тебиквюэ элььэефэнд мэдиокретатым &
            Чэнзэрет мныжаркхюм         \\
            \midrule %%% тонкий разделитель. Отделяет названия столбцов. Обязателен по ГОСТ 2.105 пункт 4.4.5
            Эй, жлоб! Где туз? Прячь юных съёмщиц в~шкаф Плюш изъят. Бьём чуждый цен хвощ! &
            \({\approx}\) &
            \({\approx}\) &
            \({\approx}\) &
            \( + \) \\
            Эх, чужак! Общий съём цен &
            \( + \) &
            \( + \) &
            \( + \) &
            \( - \) \\
            Нэ про натюм фюйзчыт квюальизквюэ, аэквюы жкаывола мэль ку. Ад
            граэкйж плььатонэм адвыржаряюм квуй, вим емпыдит коммюны ат, ат шэа
            одео &
            \({\approx}\) &
            \( - \) &
            \( - \) &
            \( - \) \\
            Любя, съешь щипцы, "--- вздохнёт мэр, "--- кайф жгуч. &
            \( - \) &
            \( + \) &
            \( + \) &
            \({\approx}\) \\
            Нэ про натюм фюйзчыт квюальизквюэ, аэквюы жкаывола мэль ку. Ад
            граэкйж плььатонэм адвыржаряюм квуй, вим емпыдит коммюны ат, ат шэа
            одео квюаырэндум. Вёртюты ажжынтиор эффикеэнди эож нэ. &
            \( + \) &
            \( - \) &
            \({\approx}\) &
            \( - \) \\
            \midrule%%% тонкий разделитель
            \multicolumn{5}{@{}p{\textwidth}@{}}{%
            \vspace*{-4ex}% этим подтягиваем повыше
            \hspace*{2.5em}% абзацный отступ - требование ГОСТ 2.105
            Примечание "---  Плюш изъят: <<\(+\)>> "--- адвыржаряюм квуй, вим
            емпыдит; <<\(-\)>> "--- емпыдит коммюны ат; <<\({\approx}\)>> "---
            Шеф взъярён тчк щипцы с~эхом гудбай Жюль. Эй, жлоб! Где туз?
            Прячь юных съёмщиц в~шкаф. Экс-граф?
            }
            \\
            \bottomrule %%% нижняя линейка
        \end{tabulary}%
    \end{SingleSpace}
\end{table}

Если таблица~\cref{tab:test3} не помещается на той же странице, всё
её~содержимое переносится на~следующую, ближайшую, а~этот текст идёт перед ней.
\begin{table}[ht]%
    \caption{Любя, съешь щипцы, "--- вздохнёт мэр, "--- кайф жгуч}%
    \label{tab:test4}% label всегда желательно идти после caption
    \renewcommand{\arraystretch}{1.6}%% Увеличение расстояния между рядами, для улучшения восприятия.
    \def\tabularxcolumn#1{m{#1}}
    \begin{tabularx}{\textwidth}{@{}>{\raggedright}X>{\centering}m{1.9cm} >{\centering}m{1.9cm} >{\centering}m{1.9cm} >{\centering\arraybackslash}m{1.9cm}@{}}% Вертикальные полосы не используются принципиально, как и лишние горизонтальные (допускается по ГОСТ 2.105 пункт 4.4.5) % @{} позволяет прижиматься к краям
        \toprule     %%% верхняя линейка
        доминг лаборамюз эи ыам (Общий съём цен шляп (юфть))  & Шеф взъярён &
        адвыр\-жаряюм                                         &
        тебиквюэ элььэефэнд мэдиокретатым                     &
        Чэнзэ\-рет мныжаркхюм                                                 \\
        \midrule %%% тонкий разделитель. Отделяет названия столбцов. Обязателен по ГОСТ 2.105 пункт 4.4.5
        Эй, жлоб! Где туз? Прячь юных съёмщиц в~шкаф Плюш изъят.
        Бьём чуждый цен хвощ!                                 &
        \({\approx}\)                                         &
        \({\approx}\)                                         &
        \({\approx}\)                                         &
        \( + \)                                                               \\
        Эх, чужак! Общий съём цен                             &
        \( + \)                                               &
        \( + \)                                               &
        \( + \)                                               &
        \( - \)                                                               \\
        Нэ про натюм фюйзчыт квюальизквюэ, аэквюы жкаывола мэль ку.
        Ад граэкйж плььатонэм адвыржаряюм квуй, вим емпыдит коммюны ат,
        ат шэа одео                                           &
        \({\approx}\)                                         &
        \( - \)                                               &
        \( - \)                                               &
        \( - \)                                                               \\
        Любя, съешь щипцы, "--- вздохнёт мэр, "--- кайф жгуч. &
        \( - \)                                               &
        \( + \)                                               &
        \( + \)                                               &
        \({\approx}\)                                                         \\
        Нэ про натюм фюйзчыт квюальизквюэ, аэквюы жкаывола мэль ку. Ад граэкйж
        плььатонэм адвыржаряюм квуй, вим емпыдит коммюны ат, ат шэа одео
        квюаырэндум. Вёртюты ажжынтиор эффикеэнди эож нэ.     &
        \( + \)                                               &
        \( - \)                                               &
        \({\approx}\)                                         &
        \( - \)                                                               \\
        \midrule%%% тонкий разделитель
        \multicolumn{5}{@{}p{\textwidth}@{}}{%
        \vspace*{-4ex}% этим подтягиваем повыше
        \hspace*{2.5em}% абзацный отступ - требование ГОСТ 2.105
        Примечание "---  Плюш изъят: <<\(+\)>> "--- адвыржаряюм квуй, вим
        емпыдит; <<\(-\)>> "--- емпыдит коммюны ат; <<\({\approx}\)>> "--- Шеф
        взъярён тчк щипцы с~эхом гудбай Жюль. Эй, жлоб! Где туз? Прячь юных
        съёмщиц в~шкаф. Экс-граф?
        }
        \\
        \bottomrule %%% нижняя линейка
    \end{tabularx}%
\end{table}

\section{Таблицы с форматированными числами}\label{sec:ch5/formatted-numbers}

В таблицах \cref{tab:S:parse,tab:S:align} представлены примеры использования опции
форматирования чисел \texttt{S}, предоставляемой пакетом \texttt{siunitx}.

\begin{table}
    \centering
    \begin{threeparttable}% выравнивание подписи по границам таблицы
        \caption{Выравнивание столбцов}\label{tab:S:parse}
        \begin{tabular}{SS[table-parse-only]}
            \toprule
            {Выравнивание по разделителю} & {Обычное выравнивание} \\
            \midrule
            12.345                        & 12.345                 \\
            6,78                          & 6,78                   \\
            -88.8(9)                      & -88.8(9)               \\
            4.5e3                         & 4.5e3                  \\
            \bottomrule
        \end{tabular}
    \end{threeparttable}
\end{table}

\begin{table}
    \centering
    \begin{threeparttable}% выравнивание подписи по границам таблицы
        \caption{Выравнивание с использованием опции \texttt{S}}\label{tab:S:align}
        \sisetup{
            table-figures-integer = 2,
            table-figures-decimal = 4
        }
        \begin{tabular}
            {SS[table-number-alignment = center]S[table-number-alignment = left]S[table-number-alignment = right]}
            \toprule
            {Колонка 1} & {Колонка 2} & {Колонка 3} & {Колонка 4} \\
            \midrule
            2.3456      & 2.3456      & 2.3456      & 2.3456      \\
            34.2345     & 34.2345     & 34.2345     & 34.2345     \\
            56.7835     & 56.7835     & 56.7835     & 56.7835     \\
            90.473      & 90.473      & 90.473      & 90.473      \\
            \bottomrule
        \end{tabular}
    \end{threeparttable}
\end{table}

\section{Параграф \texorpdfstring{\cyrdash{}}{---} два}\label{sec:ch5/sect2}
% Не все (xe|lua)latex совместимые шрифты умеют работать с русским тире "---

Некоторый текст.

\section{Параграф с подпараграфами}\label{sec:ch5/sect3}

\subsection{Подпараграф \texorpdfstring{\cyrdash{}}{---} один}\label{subsec:ch5/sect3/sub1}

Некоторый текст.

\subsection{Подпараграф \texorpdfstring{\cyrdash{}}{---} два}\label{subsec:ch5/sect3/sub2}

Некоторый текст.

Для крупных листингов есть два способа. Первый красивый, но в нём могут быть
проблемы с поддержкой кириллицы (у вас может встречаться в~комментариях
и~печатаемых сообщениях), он представлен на листинге~\cref{lst:hwbeauty}.
\begin{ListingEnv}[!h]% настройки floating аналогичны окружению figure
    \captiondelim{ } % разделитель идентификатора с номером от наименования
    \caption{Программа ,,Hello, world`` на \protect\cpp}\label{lst:hwbeauty}
    % окружение учитывает пробелы и табуляции и применяет их в сответсвии с настройками
    \begin{lstlisting}[language={[ISO]C++}]
	#include <iostream>
	using namespace std;

	int main() //кириллица в комментариях при xelatex и lualatex имеет проблемы с пробелами
	{
		cout << "Hello, world" << endl; //latin letters in commentaries
		system("pause");
		return 0;
	}
    \end{lstlisting}
\end{ListingEnv}%
Второй не~такой красивый, но без ограничений (см.~листинг~\cref{lst:hwplain}).
\begin{ListingEnv}[!h]
    \captiondelim{ } % разделитель идентификатора с номером от наименования
    \caption{Программа ,,Hello, world`` без подсветки}\label{lst:hwplain}
    \begin{Verb}

        #include <iostream>
        using namespace std;

        int main() //кириллица в комментариях
        {
            cout << "Привет, мир" << endl;
        }
    \end{Verb}
\end{ListingEnv}

Можно использовать первый для вставки небольших фрагментов
внутри текста, а второй для вставки полного
кода в приложении, если таковое имеется.

Если нужно вставить совсем короткий пример кода (одна или две строки),
то~выделение  линейками и нумерация может смотреться чересчур громоздко.
В таких случаях можно использовать окружения \texttt{lstlisting} или
\texttt{Verb} без \texttt{ListingEnv}. Приведём такой пример
с указанием языка программирования, отличного от~заданного по умолчанию:
\begin{lstlisting}[language=Haskell]
fibs = 0 : 1 : zipWith (+) fibs (tail fibs)
\end{lstlisting}
Такое решение "--- со вставкой нумерованных листингов покрупнее
и~вставок без выделения для маленьких фрагментов "--- выбрано,
например, в~книге Эндрю Таненбаума и Тодда Остина по архитектуре
компьютера.

Наконец, для оформления идентификаторов внутри строк
(функция \lstinline{main} и~тому подобное) используется
\texttt{lstinline} или, самое простое, моноширинный текст
(\texttt{\textbackslash texttt}).

Пример~\cref{lst:internal3}, иллюстрирующий подключение переопределённого
языка. Может быть полезным, если подсветка кода работает криво. Без
дополнительного окружения, с подписью и ссылкой, реализованной встроенным
средством.
\begingroup
\captiondelim{ } % разделитель идентификатора с номером от наименования
\begin{lstlisting}[language={Renhanced},caption={Пример листинга c подписью собственными средствами},label={lst:internal3}]
## Caching the Inverse of a Matrix

## Matrix inversion is usually a costly computation and there may be some
## benefit to caching the inverse of a matrix rather than compute it repeatedly
## This is a pair of functions that cache the inverse of a matrix.

## makeCacheMatrix creates a special "matrix" object that can cache its inverse

makeCacheMatrix <- function(x = matrix()) {#кириллица в комментариях при xelatex и lualatex имеет проблемы с пробелами
    i <- NULL
    set <- function(y) {
        x <<- y
        i <<- NULL
    }
    get <- function() x
    setSolved <- function(solve) i <<- solve
    getSolved <- function() i
    list(set = set, get = get,
    setSolved = setSolved,
    getSolved = getSolved)

}


## cacheSolve computes the inverse of the special "matrix" returned by
## makeCacheMatrix above. If the inverse has already been calculated (and the
## matrix has not changed), then the cachesolve should retrieve the inverse from
## the cache.

cacheSolve <- function(x, ...) {
    ## Return a matrix that is the inverse of 'x'
    i <- x$getSolved()
    if(!is.null(i)) {
        message("getting cached data")
        return(i)
    }
    data <- x$get()
    i <- solve(data, ...)
    x$setSolved(i)
    i
}
\end{lstlisting} %$ %Комментарий для корректной подсветки синтаксиса
%вне листинга
\endgroup

Листинг~\cref{lst:external1} подгружается из внешнего файла. Приходится
загружать без окружения дополнительного. Иначе по страницам не переносится.
\begingroup
\captiondelim{ } % разделитель идентификатора с номером от наименования
\lstinputlisting[lastline=78,language={R},caption={Листинг из внешнего файла},label={lst:external1}]{listings/run_analysis.R}
\endgroup

\chapter{Очень длинное название второго приложения, в~котором продемонстрирована работа с~длинными таблицами}\label{app:B}

\section{Подраздел приложения}\label{app:B1}
Вот размещается длинная таблица:
\makeatletter
\@ifpackagelater{longtable}{2024/07/04} % hotfix of bug fixed in https://github.com/latex3/latex2e/commit/7c96b6b90a730278903e71a482d88479789a89a3
{% Если много longtable* используется и новый latex, то три строки ниже правильней в главную преамбулу унести
\renewenvironment{longtable*}
  {\renewcommand\LTcaptype{}\longtable}
  {\endlongtable}
}
{\@ifpackagelater{longtable}{2024/04/26}
    {\addtocounter{table}{-1}}
    {}
}
\makeatother
\fontsize{10pt}{10pt}\selectfont
\begin{longtable*}[c]{|l|c|l|l|} %longtable* появляется из пакета ltcaption и даёт ненумерованную таблицу
    \hline
    Параметр & Умолч. & Тип & Описание               \\ \hline
    \endfirsthead   \hline
    \multicolumn{4}{|c|}{\small\slshape (продолжение)}        \\ \hline
    Параметр & Умолч. & Тип & Описание               \\ \hline
    \endhead        \hline
    \multicolumn{4}{|r|}{\small\slshape продолжение следует}  \\ \hline
    \endfoot        \hline
    \endlastfoot
    \multicolumn{4}{|l|}{\&INP}        \\ \hline
    kick & 1 & int & 0: инициализация без шума (\(p_s = const\)) \\
    &   &     & 1: генерация белого шума                  \\
    &   &     & 2: генерация белого шума симметрично относительно \\
    & & & экватора    \\
    mars & 0 & int & 1: инициализация модели для планеты Марс     \\
    kick & 1 & int & 0: инициализация без шума (\(p_s = const\)) \\
    &   &     & 1: генерация белого шума                  \\
    &   &     & 2: генерация белого шума симметрично относительно \\
    & & & экватора    \\
    mars & 0 & int & 1: инициализация модели для планеты Марс     \\
    kick & 1 & int & 0: инициализация без шума (\(p_s = const\)) \\
    &   &     & 1: генерация белого шума                  \\
    &   &     & 2: генерация белого шума симметрично относительно \\
    & & & экватора    \\
    mars & 0 & int & 1: инициализация модели для планеты Марс     \\
    kick & 1 & int & 0: инициализация без шума (\(p_s = const\)) \\
    &   &     & 1: генерация белого шума                  \\
    &   &     & 2: генерация белого шума симметрично относительно \\
    & & & экватора    \\
    mars & 0 & int & 1: инициализация модели для планеты Марс     \\
    kick & 1 & int & 0: инициализация без шума (\(p_s = const\)) \\
    &   &     & 1: генерация белого шума                  \\
    &   &     & 2: генерация белого шума симметрично относительно \\
    & & & экватора    \\
    mars & 0 & int & 1: инициализация модели для планеты Марс     \\
    kick & 1 & int & 0: инициализация без шума (\(p_s = const\)) \\
    &   &     & 1: генерация белого шума                  \\
    &   &     & 2: генерация белого шума симметрично относительно \\
    & & & экватора    \\
    mars & 0 & int & 1: инициализация модели для планеты Марс     \\
    kick & 1 & int & 0: инициализация без шума (\(p_s = const\)) \\
    &   &     & 1: генерация белого шума                  \\
    &   &     & 2: генерация белого шума симметрично относительно \\
    & & & экватора    \\
    mars & 0 & int & 1: инициализация модели для планеты Марс     \\
    kick & 1 & int & 0: инициализация без шума (\(p_s = const\)) \\
    &   &     & 1: генерация белого шума                  \\
    &   &     & 2: генерация белого шума симметрично относительно \\
    & & & экватора    \\
    mars & 0 & int & 1: инициализация модели для планеты Марс     \\
    kick & 1 & int & 0: инициализация без шума (\(p_s = const\)) \\
    &   &     & 1: генерация белого шума                  \\
    &   &     & 2: генерация белого шума симметрично относительно \\
    & & & экватора    \\
    mars & 0 & int & 1: инициализация модели для планеты Марс     \\
    kick & 1 & int & 0: инициализация без шума (\(p_s = const\)) \\
    &   &     & 1: генерация белого шума                  \\
    &   &     & 2: генерация белого шума симметрично относительно \\
    & & & экватора    \\
    mars & 0 & int & 1: инициализация модели для планеты Марс     \\
    kick & 1 & int & 0: инициализация без шума (\(p_s = const\)) \\
    &   &     & 1: генерация белого шума                  \\
    &   &     & 2: генерация белого шума симметрично относительно \\
    & & & экватора    \\
    mars & 0 & int & 1: инициализация модели для планеты Марс     \\
    kick & 1 & int & 0: инициализация без шума (\(p_s = const\)) \\
    &   &     & 1: генерация белого шума                  \\
    &   &     & 2: генерация белого шума симметрично относительно \\
    & & & экватора    \\
    mars & 0 & int & 1: инициализация модели для планеты Марс     \\
    kick & 1 & int & 0: инициализация без шума (\(p_s = const\)) \\
    &   &     & 1: генерация белого шума                  \\
    &   &     & 2: генерация белого шума симметрично относительно \\
    & & & экватора    \\
    mars & 0 & int & 1: инициализация модели для планеты Марс     \\
    kick & 1 & int & 0: инициализация без шума (\(p_s = const\)) \\
    &   &     & 1: генерация белого шума                  \\
    &   &     & 2: генерация белого шума симметрично относительно \\
    & & & экватора    \\
    mars & 0 & int & 1: инициализация модели для планеты Марс     \\
    kick & 1 & int & 0: инициализация без шума (\(p_s = const\)) \\
    &   &     & 1: генерация белого шума                  \\
    &   &     & 2: генерация белого шума симметрично относительно \\
    & & & экватора    \\
    mars & 0 & int & 1: инициализация модели для планеты Марс     \\
    \hline
    \multicolumn{4}{|l|}{\&SURFPAR}        \\ \hline
    kick & 1 & int & 0: инициализация без шума (\(p_s = const\)) \\
    &   &     & 1: генерация белого шума                  \\
    &   &     & 2: генерация белого шума симметрично относительно \\
    & & & экватора    \\
    mars & 0 & int & 1: инициализация модели для планеты Марс     \\
    kick & 1 & int & 0: инициализация без шума (\(p_s = const\)) \\
    &   &     & 1: генерация белого шума                  \\
    &   &     & 2: генерация белого шума симметрично относительно \\
    & & & экватора    \\
    mars & 0 & int & 1: инициализация модели для планеты Марс     \\
    kick & 1 & int & 0: инициализация без шума (\(p_s = const\)) \\
    &   &     & 1: генерация белого шума                  \\
    &   &     & 2: генерация белого шума симметрично относительно \\
    & & & экватора    \\
    mars & 0 & int & 1: инициализация модели для планеты Марс     \\
    kick & 1 & int & 0: инициализация без шума (\(p_s = const\)) \\
    &   &     & 1: генерация белого шума                  \\
    &   &     & 2: генерация белого шума симметрично относительно \\
    & & & экватора    \\
    mars & 0 & int & 1: инициализация модели для планеты Марс     \\
    kick & 1 & int & 0: инициализация без шума (\(p_s = const\)) \\
    &   &     & 1: генерация белого шума                  \\
    &   &     & 2: генерация белого шума симметрично относительно \\
    & & & экватора    \\
    mars & 0 & int & 1: инициализация модели для планеты Марс     \\
    kick & 1 & int & 0: инициализация без шума (\(p_s = const\)) \\
    &   &     & 1: генерация белого шума                  \\
    &   &     & 2: генерация белого шума симметрично относительно \\
    & & & экватора    \\
    mars & 0 & int & 1: инициализация модели для планеты Марс     \\
    kick & 1 & int & 0: инициализация без шума (\(p_s = const\)) \\
    &   &     & 1: генерация белого шума                  \\
    &   &     & 2: генерация белого шума симметрично относительно \\
    & & & экватора    \\
    mars & 0 & int & 1: инициализация модели для планеты Марс     \\
    kick & 1 & int & 0: инициализация без шума (\(p_s = const\)) \\
    &   &     & 1: генерация белого шума                  \\
    &   &     & 2: генерация белого шума симметрично относительно \\
    & & & экватора    \\
    mars & 0 & int & 1: инициализация модели для планеты Марс     \\
    kick & 1 & int & 0: инициализация без шума (\(p_s = const\)) \\
    &   &     & 1: генерация белого шума                  \\
    &   &     & 2: генерация белого шума симметрично относительно \\
    & & & экватора    \\
    mars & 0 & int & 1: инициализация модели для планеты Марс     \\
    \hline
\end{longtable*}

\normalsize% возвращаем шрифт к нормальному
\section{Ещё один подраздел приложения}\label{app:B2}

Нужно больше подразделов приложения!
Конвынёры витюпырата но нам, тебиквюэ мэнтётюм позтюлант ед про. Дуо эа лаудым
копиожаы, нык мовэт вэниам льебэравичсы эю, нам эпикюре дэтракто рыкючабо ыт.

Пример длинной таблицы с записью продолжения по ГОСТ 2.105:

\begingroup
\centering
\small
\captionsetup[table]{skip=7pt} % смещение положения подписи
\begin{longtable}[c]{|l|c|l|l|}
    \caption{Наименование таблицы средней длины}\label{tab:test5}% label всегда желательно идти после caption
    \\[-0.45\onelineskip]
    \hline
    Параметр & Умолч. & Тип & Описание                                          \\ \hline
    \endfirsthead%
    \caption*{Продолжение таблицы~\thetable}                                    \\[-0.45\onelineskip]
    \hline
    Параметр & Умолч. & Тип & Описание                                          \\ \hline
    \endhead
    \hline
    \endfoot
    \hline
    \endlastfoot
    \multicolumn{4}{|l|}{\&INP}                                                 \\ \hline
    kick     & 1      & int & 0: инициализация без шума (\(p_s = const\))       \\
             &        &     & 1: генерация белого шума                          \\
             &        &     & 2: генерация белого шума симметрично относительно \\
             &        &     & экватора                                          \\
    mars     & 0      & int & 1: инициализация модели для планеты Марс          \\
    kick     & 1      & int & 0: инициализация без шума (\(p_s = const\))       \\
             &        &     & 1: генерация белого шума                          \\
             &        &     & 2: генерация белого шума симметрично относительно \\
             &        &     & экватора                                          \\
    mars     & 0      & int & 1: инициализация модели для планеты Марс          \\
    kick     & 1      & int & 0: инициализация без шума (\(p_s = const\))       \\
             &        &     & 1: генерация белого шума                          \\
             &        &     & 2: генерация белого шума симметрично относительно \\
             &        &     & экватора                                          \\
    mars     & 0      & int & 1: инициализация модели для планеты Марс          \\
    kick     & 1      & int & 0: инициализация без шума (\(p_s = const\))       \\
             &        &     & 1: генерация белого шума                          \\
             &        &     & 2: генерация белого шума симметрично относительно \\
             &        &     & экватора                                          \\
    mars     & 0      & int & 1: инициализация модели для планеты Марс          \\
    kick     & 1      & int & 0: инициализация без шума (\(p_s = const\))       \\
             &        &     & 1: генерация белого шума                          \\
             &        &     & 2: генерация белого шума симметрично относительно \\
             &        &     & экватора                                          \\
    mars     & 0      & int & 1: инициализация модели для планеты Марс          \\
    kick     & 1      & int & 0: инициализация без шума (\(p_s = const\))       \\
             &        &     & 1: генерация белого шума                          \\
             &        &     & 2: генерация белого шума симметрично относительно \\
             &        &     & экватора                                          \\
    mars     & 0      & int & 1: инициализация модели для планеты Марс          \\
    kick     & 1      & int & 0: инициализация без шума (\(p_s = const\))       \\
             &        &     & 1: генерация белого шума                          \\
             &        &     & 2: генерация белого шума симметрично относительно \\
             &        &     & экватора                                          \\
    mars     & 0      & int & 1: инициализация модели для планеты Марс          \\
    kick     & 1      & int & 0: инициализация без шума (\(p_s = const\))       \\
             &        &     & 1: генерация белого шума                          \\
             &        &     & 2: генерация белого шума симметрично относительно \\
             &        &     & экватора                                          \\
    mars     & 0      & int & 1: инициализация модели для планеты Марс          \\
    kick     & 1      & int & 0: инициализация без шума (\(p_s = const\))       \\
             &        &     & 1: генерация белого шума                          \\
             &        &     & 2: генерация белого шума симметрично относительно \\
             &        &     & экватора                                          \\
    mars     & 0      & int & 1: инициализация модели для планеты Марс          \\
    kick     & 1      & int & 0: инициализация без шума (\(p_s = const\))       \\
             &        &     & 1: генерация белого шума                          \\
             &        &     & 2: генерация белого шума симметрично относительно \\
             &        &     & экватора                                          \\
    mars     & 0      & int & 1: инициализация модели для планеты Марс          \\
    kick     & 1      & int & 0: инициализация без шума (\(p_s = const\))       \\
             &        &     & 1: генерация белого шума                          \\
             &        &     & 2: генерация белого шума симметрично относительно \\
             &        &     & экватора                                          \\
    mars     & 0      & int & 1: инициализация модели для планеты Марс          \\
    kick     & 1      & int & 0: инициализация без шума (\(p_s = const\))       \\
             &        &     & 1: генерация белого шума                          \\
             &        &     & 2: генерация белого шума симметрично относительно \\
             &        &     & экватора                                          \\
    mars     & 0      & int & 1: инициализация модели для планеты Марс          \\
    kick     & 1      & int & 0: инициализация без шума (\(p_s = const\))       \\
             &        &     & 1: генерация белого шума                          \\
             &        &     & 2: генерация белого шума симметрично относительно \\
             &        &     & экватора                                          \\
    mars     & 0      & int & 1: инициализация модели для планеты Марс          \\
    kick     & 1      & int & 0: инициализация без шума (\(p_s = const\))       \\
             &        &     & 1: генерация белого шума                          \\
             &        &     & 2: генерация белого шума симметрично относительно \\
             &        &     & экватора                                          \\
    mars     & 0      & int & 1: инициализация модели для планеты Марс          \\
    kick     & 1      & int & 0: инициализация без шума (\(p_s = const\))       \\
             &        &     & 1: генерация белого шума                          \\
             &        &     & 2: генерация белого шума симметрично относительно \\
             &        &     & экватора                                          \\
    mars     & 0      & int & 1: инициализация модели для планеты Марс          \\
    \hline
    \multicolumn{4}{|l|}{\&SURFPAR}                                             \\ \hline
    kick     & 1      & int & 0: инициализация без шума (\(p_s = const\))       \\
             &        &     & 1: генерация белого шума                          \\
             &        &     & 2: генерация белого шума симметрично относительно \\
             &        &     & экватора                                          \\
    mars     & 0      & int & 1: инициализация модели для планеты Марс          \\
    kick     & 1      & int & 0: инициализация без шума (\(p_s = const\))       \\
             &        &     & 1: генерация белого шума                          \\
             &        &     & 2: генерация белого шума симметрично относительно \\
             &        &     & экватора                                          \\
    mars     & 0      & int & 1: инициализация модели для планеты Марс          \\
    kick     & 1      & int & 0: инициализация без шума (\(p_s = const\))       \\
             &        &     & 1: генерация белого шума                          \\
             &        &     & 2: генерация белого шума симметрично относительно \\
             &        &     & экватора                                          \\
    mars     & 0      & int & 1: инициализация модели для планеты Марс          \\
    kick     & 1      & int & 0: инициализация без шума (\(p_s = const\))       \\
             &        &     & 1: генерация белого шума                          \\
             &        &     & 2: генерация белого шума симметрично относительно \\
             &        &     & экватора                                          \\
    mars     & 0      & int & 1: инициализация модели для планеты Марс          \\
    kick     & 1      & int & 0: инициализация без шума (\(p_s = const\))       \\
             &        &     & 1: генерация белого шума                          \\
             &        &     & 2: генерация белого шума симметрично относительно \\
             &        &     & экватора                                          \\
    mars     & 0      & int & 1: инициализация модели для планеты Марс          \\
    kick     & 1      & int & 0: инициализация без шума (\(p_s = const\))       \\
             &        &     & 1: генерация белого шума                          \\
             &        &     & 2: генерация белого шума симметрично относительно \\
             &        &     & экватора                                          \\
    mars     & 0      & int & 1: инициализация модели для планеты Марс          \\
    kick     & 1      & int & 0: инициализация без шума (\(p_s = const\))       \\
             &        &     & 1: генерация белого шума                          \\
             &        &     & 2: генерация белого шума симметрично относительно \\
             &        &     & экватора                                          \\
    mars     & 0      & int & 1: инициализация модели для планеты Марс          \\
    kick     & 1      & int & 0: инициализация без шума (\(p_s = const\))       \\
             &        &     & 1: генерация белого шума                          \\
             &        &     & 2: генерация белого шума симметрично относительно \\
             &        &     & экватора                                          \\
    mars     & 0      & int & 1: инициализация модели для планеты Марс          \\
    kick     & 1      & int & 0: инициализация без шума (\(p_s = const\))       \\
             &        &     & 1: генерация белого шума                          \\
             &        &     & 2: генерация белого шума симметрично относительно \\
             &        &     & экватора                                          \\
    mars     & 0      & int & 1: инициализация модели для планеты Марс          \\
\end{longtable}
\normalsize% возвращаем шрифт к нормальному
\endgroup
\section{Использование длинных таблиц с окружением \textit{longtblr} из~пакета \texttt{tabularray}}\label{app:B2a}

В таблице \cref{tab:test-functions} более книжный вариант длинной таблицы,
используя окружение \verb!longtblr! из~пакета \verb!tabularray! и разнообразные
разделители (\verb!toprule!, \verb!midrule!, \verb!bottomrule!) из~пакета
\verb!booktabs!.

Чтобы визуально таблица смотрелась лучше, можно использовать следующие
параметры.
Таблица задаётся на всю ширину, \verb!longtblr! позволяет делить ширину колонок
пропорционально "--- тут три колонки в~пропорции 1.1:1.1:4 "--- для каждой
колонки первый параметр в~описании \verb!X[]!.
Кроме того, в~таблице убраны отступы слева и справа с~помощью \verb!@{}!
в~преамбуле таблицы.
К~первому и~второму столбцу применяется модификатор

\verb!>{\setlength{\baselineskip}{0.7\baselineskip}}!,

\noindent который уменьшает межстрочный интервал для текста таблиц (иначе
заголовок второго столбца значительно шире, а двухстрочное имя
сливается с~окружающими). Для первой и второй колонки текст в ячейках
выравниваются по~центру как по~вертикали, так и по горизонтали "---
задаётся буквами \verb!m!~и~\verb!c!~в~описании столбца \verb!X[]!.

Так как формулы большие "--- используется окружение \verb!alignedat!,
чтобы отступ был одинаковый у всех формул "--- он сделан для всех, хотя
для большей части можно было и не использовать.  Чтобы формулы
занимали поменьше места в~каждом столбце формулы (где надо)
используется \verb!\textstyle! "--- он~делает дроби меньше, у~знаков
суммы и произведения "--- индексы сбоку. Иногда формула слишком большая,
сливается со следующей, поэтому после неё ставится небольшой
дополнительный отступ \verb!\vspace*{2ex}!. Для штрафных функций "---
размер фигурных скобок задан вручную \verb!\Big\{!, т.\:к. не~умеет
\verb!alignedat! работать с~\verb!\left! и~\verb!\right! через
несколько строк/колонок.

В примечании к таблице наоборот, окружение \verb!cases! даёт слишком
большие промежутки между вариантами, чтобы их уменьшить, в конце
каждой строчки окружения использовался отрицательный дополнительный
отступ \verb!\\[-0.5em]!.

\DefTblrTemplate{contfoot-text}{default}{\small\slshape продолжение следует} % переделали default шаблон, используемый по-умолчанию
\DefTblrTemplate{conthead-text}{default}{\small\slshape (продолжение)} % переделали normal default, используемый по-умолчанию
\DefTblrTemplate{capcont}{default}{\centering\UseTblrTemplate{conthead-text}{default}\par} % для работы центрирования обязательно \par
\DefTblrTemplate{caplast}{default}{\small\slshape (окончание)}
\DefTblrTemplate{lasthead}{default}{\centering\UseTblrTemplate{caplast}{default}\par} % для работы центрирования обязательно \par
\DefTblrTemplate{firsthead}{default}{% правим шаблон у первого заголовка, чтобы считывал настройки из пакета caption
    % https://tex.stackexchange.com/a/628973
    \addtocounter{table}{-1}%
    \IfTokenListEmpty{\InsertTblrText{entry}}{% важно, чтобы не дублировались записи в списке таблиц
        \captionof{table}{\InsertTblrText{caption}}%
    }{%
        \captionof{table}[\InsertTblrText{entry}]{\InsertTblrText{caption}}%
    }% если будет запись в entry, то она пойдет в список таблиц, см. документацию tabularray
}
\SetTblrTemplate{caption-lot}{empty} % важно, чтобы не дублировались записи в списке таблиц
\begin{longtblr}[
    caption = {Тестовые функции для оптимизации, \(D\) "---  размерность. Для всех функций значение в точке глобального минимума равно нулю.},
    label = {tab:test-functions},
    ]{
    colspec = {%
    @{}>{\setlength{\baselineskip}{0.7\baselineskip}}X[1.1,m,c]%
    >{\setlength{\baselineskip}{0.7\baselineskip}}X[1.1,m,c]%
    X[4,l]@{}%
    },
    width = \textwidth,
    rowhead = 1,
    rows={rowsep=3pt},
    row{1}={rowsep=2pt},
    }
    \toprule     %%% верхняя линейка
    Имя                      & Стартовый диапазон параметров   & Функция                                 \\
    \midrule %%% тонкий разделитель. Отделяет названия столбцов. Обязателен по ГОСТ 2.105 пункт 4.4.5
    сфера                    & \(\left[-100,\,100\right]^D\)   &
    \(\begin{aligned}
          \textstyle f_1(x)=\sum_{i=1}^Dx_i^2
      \end{aligned}\)                                                                 \\
    Schwefel 2.22            & \(\left[-10,\,10\right]^D\)     &
    \(\begin{aligned}
          \textstyle f_2(x)=\sum_{i=1}^D|x_i|+\prod_{i=1}^D|x_i|
      \end{aligned}\)                                              \\
    Schwefel 1.2             & \(\left[-100,\,100\right]^D\)   &
    \(\begin{aligned}
          \textstyle f_3(x)=\sum_{i=1}^D\left(\sum_{j=1}^ix_j\right)^2
      \end{aligned}\)                            \\
    Schwefel 2.21            & \(\left[-100,\,100\right]^D\)   &
    \(\begin{aligned}
          \textstyle f_4(x)=\max_i\!\left\{\left|x_i\right|\right\}
      \end{aligned}\)                                           \\
    Rosenbrock               & \(\left[-30,\,30\right]^D\)     &
    \(\begin{aligned}
          \textstyle f_5(x)=
          \sum_{i=1}^{D-1}
          \left[100\!\left(x_{i+1}-x_i^2\right)^2+(x_i-1)^2\right]
      \end{aligned}\)                      \\
    ступенчатая              & \(\left[-100,\,100\right]^D\)   &
    \(\begin{aligned}
          \textstyle f_6(x)=\sum_{i=1}^D\big\lfloor x_i+0.5\big\rfloor^2
      \end{aligned}\)                                      \\
    зашумлённая квартическая & \(\left[-1.28,\,1.28\right]^D\) &
    \(\begin{aligned}
          \textstyle f_7(x)=\sum_{i=1}^Dix_i^4+rand[0,1)
      \end{aligned}\)\vspace*{2ex}                                                      \\
    Schwefel 2.26            & \(\left[-500,\,500\right]^D\)   &
    \(\begin{aligned}
          f_8(x)= & \textstyle\sum_{i=1}^D-x_i\,\sin\sqrt{|x_i|}\,+ \\
                  & \vphantom{\sum}+ D\cdot
          418.98288727243369
      \end{aligned}\)                                          \\
    Rastrigin                & \(\left[-5.12,\,5.12\right]^D\) &
    \(\begin{aligned}
          \textstyle f_9(x)=\sum_{i=1}^D\left[x_i^2-10\,\cos(2\pi x_i)+10\right]
      \end{aligned}\)\vspace*{2ex}                              \\
    Ackley                   & \(\left[-32,\,32\right]^D\)     &
    \(\begin{aligned}
          f_{10}(x)= & \textstyle -20\, \exp\!\left(
          -0.2\sqrt{\frac{1}{D}\sum_{i=1}^Dx_i^2} \right)- \\
                     & \textstyle - \exp\left(
              \frac{1}{D}\sum_{i=1}^D\cos(2\pi x_i)  \right)
          + 20 + e
      \end{aligned}\)                               \\
    Griewank                 & \(\left[-600,\,600\right]^D\)   &
    \(\begin{aligned}
          f_{11}(x)= & \textstyle \frac{1}{4000}\sum_{i=1}^{D}x_i^2 -
          \prod_{i=1}^D\cos\left(x_i/\sqrt{i}\right) +1
      \end{aligned}\) \vspace*{3ex}                                        \\
    штрафная 1               & \(\left[-50,\,50\right]^D\)     &
    \(\begin{aligned}
          f_{12}(x)= & \textstyle \frac{\pi}{D}\Big\{ 10\,\sin^2(\pi y_1) +            \\
                     & +\textstyle \sum_{i=1}^{D-1}(y_i-1)^2
          \left[1+10\,\sin^2(\pi y_{i+1})\right] +                                     \\
                     & +(y_D-1)^2 \Big\} +\textstyle\sum_{i=1}^D u(x_i,\,10,\,100,\,4)
      \end{aligned}\) \vspace*{2ex} \\
    штрафная 2               & \(\left[-50,\,50\right]^D\)     &
    \(\begin{aligned}
          f_{13}(x)= & \textstyle 0.1 \Big\{\sin^2(3\pi x_1) +            \\
                     & +\textstyle \sum_{i=1}^{D-1}(x_i-1)^2
          \left[1+\sin^2(3 \pi x_{i+1})\right] +                          \\
                     & +(x_D-1)^2\left[1+\sin^2(2\pi x_D)\right] \Big\} + \\
                     & +\textstyle\sum_{i=1}^D u(x_i,\,5,\,100,\,4)
      \end{aligned}\)              \\
    сфера                    & \(\left[-100,\,100\right]^D\)   &
    \(\begin{aligned}
          \textstyle f_1(x)=\sum_{i=1}^Dx_i^2
      \end{aligned}\)                                                                 \\
    Schwefel 2.22            & \(\left[-10,\,10\right]^D\)     &
    \(\begin{aligned}
          \textstyle f_2(x)=\sum_{i=1}^D|x_i|+\prod_{i=1}^D|x_i|
      \end{aligned}\)                                              \\
    Schwefel 1.2             & \(\left[-100,\,100\right]^D\)   &
    \(\begin{aligned}
          \textstyle f_3(x)=\sum_{i=1}^D\left(\sum_{j=1}^ix_j\right)^2
      \end{aligned}\)                            \\
    Schwefel 2.21            & \(\left[-100,\,100\right]^D\)   &
    \(\begin{aligned}
          \textstyle f_4(x)=\max_i\!\left\{\left|x_i\right|\right\}
      \end{aligned}\)                                           \\
    Rosenbrock               & \(\left[-30,\,30\right]^D\)     &
    \(\begin{aligned}
          \textstyle f_5(x)=
          \sum_{i=1}^{D-1}
          \left[100\!\left(x_{i+1}-x_i^2\right)^2+(x_i-1)^2\right]
      \end{aligned}\)                      \\
    ступенчатая              & \(\left[-100,\,100\right]^D\)   &
    \(\begin{aligned}
          \textstyle f_6(x)=\sum_{i=1}^D\big\lfloor x_i+0.5\big\rfloor^2
      \end{aligned}\)                                      \\
    зашумлённая квартическая & \(\left[-1.28,\,1.28\right]^D\) &
    \(\begin{aligned}
          \textstyle f_7(x)=\sum_{i=1}^Dix_i^4+rand[0,1)
      \end{aligned}\)\vspace*{2ex}                                                      \\
    Schwefel 2.26            & \(\left[-500,\,500\right]^D\)   &
    \(\begin{aligned}
          f_8(x)= & \textstyle\sum_{i=1}^D-x_i\,\sin\sqrt{|x_i|}\,+ \\
                  & \vphantom{\sum}+ D\cdot
          418.98288727243369
      \end{aligned}\)                                          \\
    Rastrigin                & \(\left[-5.12,\,5.12\right]^D\) &
    \(\begin{aligned}
          \textstyle f_9(x)=\sum_{i=1}^D\left[x_i^2-10\,\cos(2\pi x_i)+10\right]
      \end{aligned}\)\vspace*{2ex}                              \\
    Ackley                   & \(\left[-32,\,32\right]^D\)     &
    \(\begin{aligned}
          f_{10}(x)= & \textstyle -20\, \exp\!\left(
          -0.2\sqrt{\frac{1}{D}\sum_{i=1}^Dx_i^2} \right)- \\
                     & \textstyle - \exp\left(
              \frac{1}{D}\sum_{i=1}^D\cos(2\pi x_i)  \right)
          + 20 + e
      \end{aligned}\)                               \\
    Griewank                 & \(\left[-600,\,600\right]^D\)   &
    \(\begin{aligned}
          f_{11}(x)= & \textstyle \frac{1}{4000}\sum_{i=1}^{D}x_i^2 -
          \prod_{i=1}^D\cos\left(x_i/\sqrt{i}\right) +1
      \end{aligned}\) \vspace*{3ex}                                        \\
    штрафная 1               & \(\left[-50,\,50\right]^D\)     &
    \(\begin{aligned}
          f_{12}(x)= & \textstyle \frac{\pi}{D}\Big\{ 10\,\sin^2(\pi y_1) +            \\
                     & +\textstyle \sum_{i=1}^{D-1}(y_i-1)^2
          \left[1+10\,\sin^2(\pi y_{i+1})\right] +                                     \\
                     & +(y_D-1)^2 \Big\} +\textstyle\sum_{i=1}^D u(x_i,\,10,\,100,\,4)
      \end{aligned}\) \vspace*{2ex} \\
    штрафная 2               & \(\left[-50,\,50\right]^D\)     &
    \(\begin{aligned}
          f_{13}(x)= & \textstyle 0.1 \Big\{\sin^2(3\pi x_1) +            \\
                     & +\textstyle \sum_{i=1}^{D-1}(x_i-1)^2
          \left[1+\sin^2(3 \pi x_{i+1})\right] +                          \\
                     & +(x_D-1)^2\left[1+\sin^2(2\pi x_D)\right] \Big\} + \\
                     & +\textstyle\sum_{i=1}^D u(x_i,\,5,\,100,\,4)
      \end{aligned}\)              \\
    \midrule%%% тонкий разделитель
    \SetCell[c=3]{l,\linewidth}%
    \hspace*{2.5em}% абзацный отступ - требование ГОСТ 2.105
    Примечание "---  Для функций \(f_{12}\) и \(f_{13}\)
    используется \(y_i = 1 + \frac{1}{4}(x_i+1)\)
    и~\(u(x_i,\,a,\,k,\,m)=
    \begin{cases*}
        k(x_i-a)^m,  & \( x_i >a \)            \\[-0.5em]
        0,           & \( -a\leq x_i \leq a \) \\[-0.5em]
        k(-x_i-a)^m, & \( x_i <-a \)
    \end{cases*}
    \)                                                                                                   \\
    \bottomrule %%% нижняя линейка
\end{longtblr}

\section{Форматирование внутри таблиц}\label{app:B3}

В таблице \cref{tab:other-row} пример с чересстрочным форматированием.
Это реализовано средствами, доступными в таблицах пакета \verb+tabularray+.

В таблице \cref{tab:other-row} каждая чётная строчка (заголовок таблицы тоже
считается за~строчку) "--- синяя, нечётная "--- с наклоном и~слегка поднята
вверх.
Визуально это приводит к тому, что среднее значение и~среднеквадратичное
изменение группируются и хорошо выделяются взглядом в~таблице.
Сохраняется возможность отдельные значения в таблице выделить цветом или
шрифтом.
К~первому и~второму столбцу форматирование не применяется по~сути таблицы,
к~шестому общее форматирование не~применяется для наглядности.

\needspace{2\baselineskip}
\begin{longtblr}[
    caption = {Длинная таблица с примером чересстрочного форматирования},
    label = {tab:other-row},
    ]{
    vspan=minimal,
    stretch=0,
    rows={m,ht=0.8\baselineskip},
    columns={c},
    colspec = {%
    @{}X[0.27,l]%
    @{}X[0.7]%
    @{}X%
    @{}X%
    @{}X%
    @{}X[0.98]%
    @{}X[preto={\setlength{\baselineskip}{0.7\baselineskip}}]%
    @{}X@{}%
    },
    width = \textwidth,
    rowhead = 1,
    cell{even}{3-5,7-8} = {font=\color{blue}},
    cell{odd[3]}{3-5,7-8} = {cmd={\vspace{0.3ex}},font=\itshape},
        }
    \toprule %%% верхняя линейка
        & Итера\-ции & JADE\texttt{++}  & JADE      & jDE                    & SaDE                   & DE/rand /1/bin & PSO       \\
    \midrule %%% тонкий разделитель. Отделяет названия столбцов. Обязателен по ГОСТ 2.105 пункт 4.4.5
    f1  & 1500       & \textbf{1.8E-60} & 1.3E-54   & 2.5E-28                & 4.5E-20                & 9.8E-14        & 9.6E-42   \\\nopagebreak
        &            & (8.4E-60)        & (9.2E-54) & {\color{red}(3.5E-28)} & (6.9E-20)              & (8.4E-14)      & (2.7E-41) \\
    f2  & 2000       & 1.8E-25          & 3.9E-22   & 1.5E-23                & 1.9E-14                & 1.6E-09        & 9.3E-21   \\\nopagebreak
        &            & (8.8E-25)        & (2.7E-21) & (1.0E-23)              & (1.1E-14)              & (1.1E-09)      & (6.3E-20) \\
    f3  & 5000       & 5.7E-61          & 6.0E-87   & 5.2E-14                & {\color{green}9.0E-37} & 6.6E-11        & 2.5E-19   \\\nopagebreak
        &            & (2.7E-60)        & (1.9E-86) & (1.1E-13)              & (5.4E-36)              & (8.8E-11)      & (3.9E-19) \\
    f4  & 5000       & 8.2E-24          & 4.3E-66   & 1.4E-15                & 7.4E-11                & 4.2E-01        & 4.4E-14   \\\nopagebreak
        &            & (4.0E-23)        & (1.2E-65) & (1.0E-15)              & (1.8E-10)              & (1.1E+00)      & (9.3E-14) \\
    f5  & 3000       & 8.0E-02          & 3.2E-01   & 1.3E+01                & 2.1E+01                & 2.1E+00        & 2.5E+01   \\\nopagebreak
        &            & (5.6E-01)        & (1.1E+00) & (1.4E+01)              & (7.8E+00)              & (1.5E+00)      & (3.2E+01) \\
    f6  & 100        & 2.9E+00          & 5.6E+00   & 1.0E+03                & 9.3E+02                & 4.7E+03        & 4.5E+01   \\\nopagebreak
        &            & (1.2E+00)        & (1.6E+00) & (2.2E+02)              & (1.8E+02)              & (1.1E+03)      & (2.4E+01) \\
    f7  & 3000       & 6.4E-04          & 6.8E-04   & 3.3E-03                & 4.8E-03                & 4.7E-03        & 2.5E-03   \\\nopagebreak
        &            & (2.5E-04)        & (2.5E-04) & (8.5E-04)              & (1.2E-03)              & (1.2E-03)      & (1.4E-03) \\
    f8  & 1000       & 3.3E-05          & 7.1E+00   & 7.9E-11                & 4.7E+00                & 5.9E+03        & 2.4E+03   \\\nopagebreak
        &            & (2.3E-05)        & (2.8E+01) & (1.3E-10)              & (3.3E+01)              & (1.1E+03)      & (6.7E+02) \\
    f9  & 1000       & 1.0E-04          & 1.4E-04   & 1.5E-04                & 1.2E-03                & 1.8E+02        & 5.2E+01   \\\nopagebreak
        &            & (6.0E-05)        & (6.5E-05) & (2.0E-04)              & (6.5E-04)              & (1.3E+01)      & (1.6E+01) \\
    f10 & 500        & 8.2E-10          & 3.0E-09   & 3.5E-04                & 2.7E-03                & 1.1E-01        & 4.6E-01   \\\nopagebreak
        &            & (6.9E-10)        & (2.2E-09) & (1.0E-04)              & (5.1E-04)              & (3.9E-02)      & (6.6E-01) \\
    f11 & 500        & 9.9E-08          & 2.0E-04   & 1.9E-05                & 7.8E-04                & 2.0E-01        & 1.3E-02   \\\nopagebreak
        &            & (6.0E-07)        & (1.4E-03) & (5.8E-05)              & (1.2E-03)              & (1.1E-01)      & (1.7E-02) \\
    f12 & 500        & 4.6E-17          & 3.8E-16   & 1.6E-07                & 1.9E-05                & 1.2E-02        & 1.9E-01   \\\nopagebreak
        &            & (1.9E-16)        & (8.3E-16) & (1.5E-07)              & (9.2E-06)              & (1.0E-02)      & (3.9E-01) \\
    f13 & 500        & 2.0E-16          & 1.2E-15   & 1.5E-06                & 6.1E-05                & 7.5E-02        & 2.9E-03   \\\nopagebreak
        &            & (6.5E-16)        & (2.8E-15) & (9.8E-07)              & (2.0E-05)              & (3.8E-02)      & (4.8E-03) \\
    f1  & 1500       & \textbf{1.8E-60} & 1.3E-54   & 2.5E-28                & 4.5E-20                & 9.8E-14        & 9.6E-42   \\\nopagebreak
        &            & (8.4E-60)        & (9.2E-54) & {\color{red}(3.5E-28)} & (6.9E-20)              & (8.4E-14)      & (2.7E-41) \\
    f2  & 2000       & 1.8E-25          & 3.9E-22   & 1.5E-23                & 1.9E-14                & 1.6E-09        & 9.3E-21   \\\nopagebreak
        &            & (8.8E-25)        & (2.7E-21) & (1.0E-23)              & (1.1E-14)              & (1.1E-09)      & (6.3E-20) \\
    f3  & 5000       & 5.7E-61          & 6.0E-87   & 5.2E-14                & 9.0E-37                & 6.6E-11        & 2.5E-19   \\\nopagebreak
        &            & (2.7E-60)        & (1.9E-86) & (1.1E-13)              & (5.4E-36)              & (8.8E-11)      & (3.9E-19) \\
    f4  & 5000       & 8.2E-24          & 4.3E-66   & 1.4E-15                & 7.4E-11                & 4.2E-01        & 4.4E-14   \\\nopagebreak
        &            & (4.0E-23)        & (1.2E-65) & (1.0E-15)              & (1.8E-10)              & (1.1E+00)      & (9.3E-14) \\
    f5  & 3000       & 8.0E-02          & 3.2E-01   & 1.3E+01                & 2.1E+01                & 2.1E+00        & 2.5E+01   \\\nopagebreak
        &            & (5.6E-01)        & (1.1E+00) & (1.4E+01)              & (7.8E+00)              & (1.5E+00)      & (3.2E+01) \\
    f6  & 100        & 2.9E+00          & 5.6E+00   & 1.0E+03                & 9.3E+02                & 4.7E+03        & 4.5E+01   \\\nopagebreak
        &            & (1.2E+00)        & (1.6E+00) & (2.2E+02)              & (1.8E+02)              & (1.1E+03)      & (2.4E+01) \\
    f7  & 3000       & 6.4E-04          & 6.8E-04   & 3.3E-03                & 4.8E-03                & 4.7E-03        & 2.5E-03   \\\nopagebreak
        &            & (2.5E-04)        & (2.5E-04) & (8.5E-04)              & (1.2E-03)              & (1.2E-03)      & (1.4E-03) \\
    f8  & 1000       & 3.3E-05          & 7.1E+00   & 7.9E-11                & 4.7E+00                & 5.9E+03        & 2.4E+03   \\\nopagebreak
        &            & (2.3E-05)        & (2.8E+01) & (1.3E-10)              & (3.3E+01)              & (1.1E+03)      & (6.7E+02) \\
    f9  & 1000       & 1.0E-04          & 1.4E-04   & 1.5E-04                & 1.2E-03                & 1.8E+02        & 5.2E+01   \\\nopagebreak
        &            & (6.0E-05)        & (6.5E-05) & (2.0E-04)              & (6.5E-04)              & (1.3E+01)      & (1.6E+01) \\
    f10 & 500        & 8.2E-10          & 3.0E-09   & 3.5E-04                & 2.7E-03                & 1.1E-01        & 4.6E-01   \\\nopagebreak
        &            & (6.9E-10)        & (2.2E-09) & (1.0E-04)              & (5.1E-04)              & (3.9E-02)      & (6.6E-01) \\
    f11 & 500        & 9.9E-08          & 2.0E-04   & 1.9E-05                & 7.8E-04                & 2.0E-01        & 1.3E-02   \\\nopagebreak
        &            & (6.0E-07)        & (1.4E-03) & (5.8E-05)              & (1.2E-03)              & (1.1E-01)      & (1.7E-02) \\
    f12 & 500        & 4.6E-17          & 3.8E-16   & 1.6E-07                & 1.9E-05                & 1.2E-02        & 1.9E-01   \\\nopagebreak
        &            & (1.9E-16)        & (8.3E-16) & (1.5E-07)              & (9.2E-06)              & (1.0E-02)      & (3.9E-01) \\
    f13 & 500        & 2.0E-16          & 1.2E-15   & 1.5E-06                & 6.1E-05                & 7.5E-02        & 2.9E-03   \\\nopagebreak
        &            & (6.5E-16)        & (2.8E-15) & (9.8E-07)              & (2.0E-05)              & (3.8E-02)      & (4.8E-03) \\
    \bottomrule %%% нижняя линейка
\end{longtblr}

\section{Стандартные префиксы ссылок}\label{app:B4}

Общепринятым является следующий формат ссылок: \texttt{<prefix>:<label>}.
Например, \verb+\label{fig:knuth}+; \verb+\ref{tab:test1}+; \verb+label={lst:external1}+.
В~таблице \cref{tab:tab_pref} приведены стандартные префиксы для различных
типов ссылок.

\begin{table}[htbp]
    \captionsetup{justification=centering}
    \centering{
        \caption{\label{tab:tab_pref}Стандартные префиксы ссылок}
        \begin{tabular}{ll}
            \toprule
            \textbf{Префикс} & \textbf{Описание} \\
            \midrule
            ch:              & Глава             \\
            sec:             & Секция            \\
            subsec:          & Подсекция         \\
            fig:             & Рисунок           \\
            tab:             & Таблица           \\
            eq:              & Уравнение         \\
            lst:             & Листинг программы \\
            itm:             & Элемент списка    \\
            alg:             & Алгоритм          \\
            app:             & Секция приложения \\
            \bottomrule
        \end{tabular}
    }
\end{table}


Для упорядочивания ссылок можно использовать разделительные символы.
Например, \verb+\label{fig:scheemes/my_scheeme}+ или \\ \verb+\label{lst:dts/linked_list}+.

\section{Очередной подраздел приложения}\label{app:B5}

Нужно больше подразделов приложения!

\section{И ещё один подраздел приложения}\label{app:B6}

Нужно больше подразделов приложения!

\clearpage
\refstepcounter{chapter}
\addcontentsline{toc}{appendix}{\protect\chapternumberline{\thechapter}Чертёж детали}

\includepdf[pages=-]{Dissertation/images/drawing.pdf}


Обзор, введение в тему, обозначение места данной работы в
мировых исследованиях и~т.\:п., можно использовать ссылки на~другие
работы~\autocite{Gosele1999161,Lermontov}
(если их~нет, то~в~автореферате
автоматически пропадёт раздел <<Список литературы>>). Внимание! Ссылки
на~другие работы в~разделе общей характеристики работы можно
использовать только при использовании \verb!biblatex! (из-за технических
ограничений \verb!bibtex8!. Это связано с тем, что одна
и~та~же~характеристика используются и~в~тексте диссертации, и в
автореферате. В~последнем, согласно ГОСТ, должен присутствовать список
работ автора по~теме диссертации, а~\verb!bibtex8! не~умеет выводить в~одном
файле два списка литературы).
При использовании \verb!biblatex! возможно использование исключительно
в~автореферате подстрочных ссылок
для других работ командой \verb!\autocite!~\autocite{Marketing}, а~также цитирование
собственных работ командой \verb!\cite!. Для этого в~файле
\verb!common/setup.tex! необходимо присвоить положительное значение
счётчику \verb!\setcounter{usefootcite}{1}!.

Для генерации содержимого титульного листа автореферата, диссертации
и~презентации используются данные из файла \verb!common/data.tex!. Если,
например, вы меняете название диссертации, то оно автоматически
появится в~итоговых файлах после очередного запуска \LaTeX. Согласно
ГОСТ 7.0.11-2011 <<5.1.1 Титульный лист является первой страницей
диссертации, служит источником информации, необходимой для обработки и
поиска документа>>. Наличие логотипа организации на~титульном листе
упрощает обработку и~поиск, для этого разметите логотип вашей
организации в папке images в~формате PDF (лучше найти его в векторном
варианте, чтобы он хорошо смотрелся при печати) под именем
\verb!logo.pdf!. Настроить размер изображения с логотипом можно
в~соответствующих местах файлов \verb!title.tex!  отдельно для
диссертации и автореферата. Если вам логотип не~нужен, то просто
удалите файл с~логотипом.

При использовании пакета \verb!biblatex! будут подсчитаны все работы, добавленные
в файл \verb!biblio/author.bib!. Для правильного подсчёта работ в~различных
системах цитирования требуется использовать поля:
\begin{itemize}
        \item \texttt{authorvak} если публикация индексирована ВАК,
        \item \texttt{authorscopus} если публикация индексирована Scopus,
        \item \texttt{authorwos} если публикация индексирована Web of Science,
        \item \texttt{authorconf} для докладов конференций,
        \item \texttt{authorpatent} для патентов,
        \item \texttt{authorprogram} для зарегистрированных программ для ЭВМ,
        \item \texttt{authorother} для других публикаций.
\end{itemize}
Для подсчёта используются счётчики:
\begin{itemize}
        \item \texttt{citeauthorvak} для работ, индексируемых ВАК,
        \item \texttt{citeauthorscopus} для работ, индексируемых Scopus,
        \item \texttt{citeauthorwos} для работ, индексируемых Web of Science,
        \item \texttt{citeauthorvakscopuswos} для работ, индексируемых одной из трёх баз,
        \item \texttt{citeauthorscopuswos} для работ, индексируемых Scopus или Web of~Science,
        \item \texttt{citeauthorconf} для докладов на конференциях,
        \item \texttt{citeauthorother} для остальных работ,
        \item \texttt{citeauthorpatent} для патентов,
        \item \texttt{citeauthorprogram} для зарегистрированных программ для ЭВМ,
        \item \texttt{citeauthor} для суммарного количества работ.
\end{itemize}
% Счётчик \texttt{citeexternal} используется для подсчёта процитированных публикаций;
% \texttt{citeregistered} "--- для подсчёта суммарного количества патентов и программ для ЭВМ.

Для добавления в список публикаций автора работ, которые не были процитированы в
автореферате, требуется их~перечислить с использованием команды \verb!\nocite! в
\verb!Synopsis/content.tex!.